% Vietnamese translation for OI Public Library
% Source-PDF: IOI中国国家候选队论文/国家集训队2013论文集.pdf
\documentclass[11pt]{article}
\usepackage{oipl-vn}
\usepackage{array}

\title{Tập luận văn đội tuyển quốc gia Trung Quốc năm 2013}
\author{Huấn luyện viên: Hu Weidong, Tang Wenbin\\Biên tập: Hu Weidong}
\date{China Computer Federation, 2013}

\begin{document}
\maketitle

\origtitle{Tập luận văn ứng viên đội tuyển quốc gia Trung Quốc Olympic Tin học năm 2013}
\sourcepath{IOI China candidate team papers / National training team 2013 collection.pdf}

\renewcommand{\abstractname}{Tóm tắt}
\begin{abstract}
PDF nguồn là tập hợp luận văn đội tuyển quốc gia Trung Quốc năm 2013. Bản
dịch này chuyển ngữ thông tin đầu sách, mục lục và sáu bài đầu tiên:
\emph{Kế hoạch lên đỉnh}, của Peng Tianyi, \emph{Lấy số trên lưới}, của
Wang Kangning, \emph{Thảo luận bài Two strings}, của Luo Gan, và
\emph{Catch The Penguins}, của Zhang Wentao, và \emph{Bàn về ứng dụng tư
tưởng chia khối trong một lớp bài xử lý dữ liệu}, của Luo Jianqiao, và
\emph{Kỹ thuật meet in the middle trong bài toán tìm kiếm}, của Qiao Mingda.
Trạng thái hiện tại: sáu bài đầu đã được dịch; các bài còn lại mới có mục
lục. Lớp văn bản của phần thân bị lỗi ánh xạ phông chữ, nên phần bài viết
được đối chiếu từ OCR trang render và các công thức, ví dụ còn đọc được
trong lớp văn bản.
\end{abstract}

\section{Thông tin đầu sách}

Tập sách có tiêu đề \emph{Tập luận văn ứng viên đội tuyển quốc gia Trung Quốc
Olympic Tin học năm 2013}. Huấn luyện viên là Hu Weidong và Tang Wenbin; biên
tập là Hu Weidong. Đơn vị xuất bản ghi trên bìa là China Computer Federation.

\section{Mục lục đã dịch}

\begin{center}
\begin{tabular}{r >{\raggedright\arraybackslash}p{0.47\linewidth} >{\raggedright\arraybackslash}p{0.28\linewidth}}
\textbf{Trang} & \textbf{Bài viết} & \textbf{Tác giả / trường} \\
\hline
1 & Kế hoạch lên đỉnh & Peng Tianyi, Trường THPT trực thuộc Đại học Sư phạm Hồ Nam \\
9 & Lấy số trên lưới & Wang Kangning, Trường THPT trực thuộc Đại học Sư phạm Đông Bắc, Cát Lâm \\
17 & Thảo luận bài \emph{Two strings} & Luo Gan, Trường Học Quân Hàng Châu, Chiết Giang \\
23 & \emph{Catch The Penguins} & Zhang Wentao, Trường Học Quân Hàng Châu, Chiết Giang \\
27 & Bàn về ứng dụng tư tưởng chia khối trong một lớp bài xử lý dữ liệu & Luo Jianqiao, Trường THPT số 80 Bắc Kinh \\
43 & Kỹ thuật \emph{meet in the middle} trong bài toán tìm kiếm & Qiao Mingda, Trường Ngoại ngữ Nam Kinh, Giang Tô \\
57 & Phân tích cơ sở và ứng dụng của xác suất trong thi tin học & Hu Yuanming, Trường Trung học Dương Châu, Giang Tô \\
71 & Bàn về một số cách giải không kinh điển cho bài cấu trúc dữ liệu & Xu Haoran, Trường Ngoại ngữ Nam Kinh, Giang Tô \\
85 & Ứng dụng cây cân bằng theo trọng lượng và cây cân bằng hậu tố trong Olympic Tin học & Chen Lijie, Trường Ngoại ngữ Hàng Châu, Chiết Giang \\
99 & Bàn về bài toán đếm trên vòng & Gao Shenghan, Trường Trung học cao cấp Thường Châu, Giang Tô \\
109 & Ứng dụng của phương pháp chia khối & Wang Ziyu, Trường Trung học số 1 Shengli, Sơn Đông \\
121 & Bàn về nguyên lý bao hàm--loại trừ & Wang Di, Trường Trung học số 7 Thành Đô, Tứ Xuyên
\end{tabular}
\end{center}

\section{Trạng thái và ghi chú trích xuất}

Các trang in trong mục lục lệch \(4\) trang so với trang vật lý của PDF: trang
vật lý \(5\) là trang in \(1\). Bài \emph{Kế hoạch lên đỉnh} dùng các trang
in \(1\)--\(8\), tức trang vật lý \(5\)--\(12\). Bài \emph{Lấy số trên lưới}
dùng các trang in \(9\)--\(16\), tức trang vật lý \(13\)--\(20\). Bài
\emph{Thảo luận bài Two strings} bắt đầu ở trang in \(17\), tức trang vật lý
\(21\); bài kế tiếp bắt đầu ở trang in \(23\), tức trang vật lý \(27\). Trang
in \(22\) là trang trắng, nên phần dịch bài thứ ba lấy nội dung ở các trang
in \(17\)--\(21\). Bài \emph{Catch The Penguins} dùng các trang in
\(23\)--\(26\), tức trang vật lý \(27\)--\(30\). Bài \emph{Bàn về ứng dụng
tư tưởng chia khối trong một lớp bài xử lý dữ liệu} bắt đầu ở trang in
\(27\), tức trang vật lý \(31\); bài kế tiếp bắt đầu ở trang in \(43\), tức
trang vật lý \(47\). Trang in \(42\) là trang trắng, nên phần dịch bài thứ
năm lấy nội dung ở các trang in \(27\)--\(41\).
Bài \emph{Kỹ thuật meet in the middle trong bài toán tìm kiếm} bắt đầu ở
trang in \(43\), tức trang vật lý \(47\); bài kế tiếp bắt đầu ở trang in
\(57\), tức trang vật lý \(61\). Phần dịch bài thứ sáu lấy nội dung ở các
trang in \(43\)--\(56\).

\section{Kế hoạch lên đỉnh}

\begin{center}
Peng Tianyi\\
Trường THPT trực thuộc Đại học Sư phạm Hồ Nam
\end{center}

\subsection{Bối cảnh bài toán}

Leo núi là một môn vận động rất có ích cho sức khỏe thể chất và tinh thần.
Nhưng khi leo núi, ta nên chọn chiến lược như thế nào để lên tới đỉnh núi
nhanh nhất? Từ câu hỏi thú vị này, tác giả nghĩ ra bài toán dưới đây.

\subsection{Mô tả bài toán}

\paragraph{Giới hạn tài nguyên.}
Giới hạn thời gian: \(2\) giây. Giới hạn bộ nhớ: \(512\) MB.

\paragraph{Phát biểu.}
Trên mặt phẳng hai chiều, một dãy núi được xác định bởi một dãy đỉnh
\[
  (x_1,y_1),(x_2,y_2),\ldots,(x_n,y_n).
\]
Hai đỉnh kề nhau được nối bằng đoạn thẳng, với bảo đảm \(x_i\) tăng nghiêm
ngặt. Như vậy ta thu được một dãy núi liên tục; giá trị \(y\) của mỗi điểm là
độ cao của điểm đó.

Ta định nghĩa phần trong của dãy núi là phần kẹp giữa đường gấp khúc nối các
đỉnh và trục \(x\), không tính chính đường gấp khúc đó. Nếu đoạn thẳng nối
đỉnh \(A\) và đỉnh \(B\) không đi qua phần trong của dãy núi, ta nói đỉnh
\(A\) nhìn thấy đỉnh \(B\), hoặc đỉnh \(B\) nhìn thấy đỉnh \(A\).

Bây giờ pty xuất phát từ một đỉnh nào đó và muốn lên tới đỉnh cao nhất của
dãy núi. Anh chỉ có thể đi trên đường gấp khúc giữa các đỉnh. Sau khi suy
nghĩ, anh chọn cách leo như sau:
\begin{enumerate}
  \item Đứng tại điểm xuất phát và nhìn sang hai bên. Nếu đỉnh núi cao nhất
  nhìn thấy được lúc này nằm bên trái, anh đi sang trái; ngược lại anh đi
  sang phải.
  \item Trong lúc đi, pty vẫn quan sát đỉnh cao nhất nhìn thấy được ở hai
  bên. Một khi phát hiện một đỉnh cao hơn tất cả những đỉnh đã thấy trước
  đó, anh sẽ đi về phía đỉnh cao nhất lúc này.
  \item Nếu có một thời điểm mà vị trí pty đang đứng cao hơn mọi đỉnh cao
  nhất nhìn thấy được ở hai bên, thì anh đã tới đỉnh cao nhất của dãy núi và
  quá trình leo núi kết thúc.
\end{enumerate}

pty muốn biết: nếu dùng chiến lược trên và xuất phát từ từng đỉnh, quãng
đường để tới đỉnh cao nhất lần lượt là bao nhiêu? Khoảng cách giữa hai điểm
trên mặt phẳng là độ dài đoạn thẳng nối chúng.

\paragraph{Dữ liệu vào.}
Dòng đầu tiên chứa một số nguyên \(n\), là số đỉnh của dãy núi.

Tiếp theo là \(n\) dòng; dòng thứ \(i\) chứa hai số nguyên \(x_i,y_i\), biểu
diễn tọa độ của đỉnh thứ \(i\). Bảo đảm \(x_i\) tăng nghiêm ngặt, các
\(y_i\) đôi một khác nhau trừ \(y_1,y_n\), và \(x_i,y_i\) đều là số nguyên
không âm. Bảo đảm \(y_1=y_n=0\).

\paragraph{Dữ liệu ra.}
In ra \(n\) dòng, mỗi dòng một số thực.

Số thực ở dòng thứ \(i\) biểu diễn quãng đường từ đỉnh thứ \(i\) tới đỉnh cao
nhất. Nếu sai số giữa kết quả in ra và đáp án chuẩn không vượt quá
\(10^{-2}\), test đó được điểm tối đa; ngược lại được \(0\) điểm.

\paragraph{Ví dụ vào.}
\begin{verbatim}
6
0 0
1 6
2 4
3 1
5 5
6 0
\end{verbatim}

\paragraph{Ví dụ ra.}
\begin{verbatim}
6.08
0.00
2.24
6.52
9.87
14.97
\end{verbatim}

\paragraph{Giải thích ví dụ.}
Các lộ trình là:
\[
\begin{array}{ll}
A: & A\to B,\\
B: & B,\\
C: & C\to B,\\
D: & D\to G\to D\to C\to B,\\
E: & E\to D\to C\to B,\\
F: & F\to E\to D\to C\to B.
\end{array}
\]
Khi xuất phát từ \(D\), đỉnh cao nhất nhìn thấy được ban đầu là \(E\). Khi đi
tới điểm \(G\), pty phát hiện đỉnh cao hơn là \(B\), nên quay đầu và tiếp tục
đi về phía \(B\). Với các điểm xuất phát còn lại, lộ trình không cần đổi
hướng.

\paragraph{Quy mô dữ liệu và quy ước.}
Mọi dữ liệu đều thỏa mãn \(x_i,y_i\le 100000\).

\begin{center}
\begin{tabular}{c l}
\textbf{Test} & \textbf{Ràng buộc} \\
\hline
1--4 & \(n\le 20\) \\
5--8 & \(n\le 70\) \\
9--10 & \(n\le 100000\), và mọi đỉnh đều nhìn thấy trực tiếp đỉnh cao nhất \\
11--14 & \(n\le 30000\) \\
15--20 & \(n\le 100000\)
\end{tabular}
\end{center}

\subsection{Phân tích bài toán}

\subsubsection{Đơn giản hóa bài toán}

Bài toán yêu cầu tại mọi thời điểm khi leo núi ta đều phải quan sát đỉnh cao
nhất hiện tại ở hai bên. Trước hết hãy tạm đơn giản hóa vấn đề: chỉ khi tới
một đỉnh ta mới quan sát đỉnh cao nhất lúc đó. Nói cách khác, điểm quan sát
tạm thời chuyển từ vô hạn nhiều điểm thành hữu hạn \(O(n)\) đỉnh.

\subsubsection{Tìm đỉnh cao nhất mà một đỉnh nhìn thấy}

Gọi đỉnh cao nhất mà đỉnh \(i\) nhìn thấy được về bên trái là \(L[i]\), và
đỉnh cao nhất mà nó nhìn thấy được về bên phải là \(R[i]\). Hãy tưởng tượng
ta đứng tại một điểm nào đó; ban đầu tia nhìn nằm ngang về bên trái, sau đó
ta từ từ ngẩng đầu lên cho tới khi tia nhìn không còn bị một ngọn núi nào che
khuất. Khi đó ta tìm được \(L[i]\).

Giả sử đỉnh hiện tại là \(A\), còn \(L[i]\) là \(B\). Hai tính chất sau là
hiển nhiên:
\begin{itemize}
  \item \(i\) điểm đầu đều nằm phía dưới đoạn \(AB\), theo nghĩa không nghiêm
  ngặt;
  \item \(A\) và \(B\) là hai điểm kề nhau cuối cùng trên bao lồi trên của
  \(i\) điểm đầu.
\end{itemize}

Vì \(A\) chắc chắn là điểm cuối cùng của bao lồi trên của \(i\) điểm đầu,
\(B\) chính là điểm áp chót trên bao lồi trên đó.\footnote{Một kết luận đẹp
như vậy thật sự làm tác giả ngạc nhiên.} Do đó ta chỉ cần dùng một stack để
duy trì bao lồi của các điểm đầu tiên. Mỗi lần thêm điểm thứ \(i+1\), ta đồng
thời duy trì tính chất của bao lồi. Vì tổng cộng chỉ có \(O(n)\) đỉnh, ta có
thể tìm toàn bộ \(L[i]\) và \(R[i]\) trong thời gian \(O(n)\).

\subsubsection{Xây dựng mô hình đồ thị}

\paragraph{Cách làm vét cạn.}
Sau khi có \(L[i]\) và \(R[i]\), ta có một thuật toán khá thô: với từng đỉnh
xuất phát, mô phỏng sự thay đổi của lộ trình. Độ phức tạp là \(O(n^2)\) hoặc
cao hơn. Điều này cho thấy ta cần tiếp tục khai thác phần thông tin bị tính
lặp.

\paragraph{Dựng cây.}
Gọi \(H[i]\) là độ cao của đỉnh thứ \(i\), và đặt
\[
  T[i]=\max(H[L[i]],H[R[i]]).
\]
Khi đó \(T[i]\) biểu diễn độ cao của đỉnh cao nhất mà điểm \(i\) nhìn thấy
được.

Khi ta xuất phát từ \(A\), đến điểm đầu tiên \(B\) thỏa mãn
\[
  T[B]\ge T[A],
\]
thì đoạn đường còn lại cần đi sẽ hoàn toàn giống như khi xuất phát từ \(B\).
Vì vậy khoảng cách từ \(A\) tới đỉnh cao nhất bằng khoảng cách từ \(B\) tới
đỉnh cao nhất cộng thêm khoảng cách từ \(A\) tới \(B\).

Lúc này cho \(B\) nối tới \(A\) bằng một cạnh có hướng, trọng số là độ dài
\(AB\). Làm như vậy với mọi đỉnh, ta thu được một cây: đỉnh cao nhất không có
bậc vào, còn mọi đỉnh khác có bậc vào bằng \(1\). Tổng trọng số trên đường đi
từ gốc cây tới một đỉnh bất kỳ chính là quãng đường từ đỉnh đó tới đỉnh cao
nhất. Chỉ cần DFS một lần là có đáp án cho mọi đỉnh.

\paragraph{Vấn đề cần giải quyết.}
Với một điểm \(A\), làm thế nào nhanh chóng tìm được điểm đầu tiên về bên
trái hoặc bên phải thỏa mãn \(T[B]\ge T[A]\)? Đây là một bài toán cấu trúc dữ
liệu rất kinh điển. Ta có thể giải bằng nhị phân đáp án kết hợp RMQ trên
đoạn. Dưới đây là một cách có hằng số nhỏ hơn.

Lập danh sách liên kết hai chiều của toàn bộ đỉnh theo thứ tự từ trái sang
phải. Sau đó duyệt tất cả các điểm theo thứ tự \(T\) tăng dần. Mỗi khi duyệt
tới một điểm \(A\), điểm đầu tiên bên trái thỏa mãn \(T[B]\ge T[A]\) chính là
hàng xóm bên trái của \(A\) còn tồn tại trong danh sách liên kết tại thời
điểm đó. Nối cạnh tương ứng, rồi xóa \(A\) khỏi danh sách liên kết. Phía bên
phải được xử lý tương tự. Độ phức tạp của thuật toán này chính là độ phức tạp
sắp xếp.

\subsubsection{Quay lại bài toán ban đầu}

Tới đây ta đã giải xong phiên bản đơn giản hóa trong thời gian sắp xếp
\(O(n\log n)\). Bây giờ hãy quay lại bài toán ban đầu.

Ta có khả năng quan sát đỉnh cao nhất tại mọi thời điểm. Điều này làm bài
toán thay đổi như thế nào?

Giả sử hiện tại ta xuất phát từ điểm \(A\), đi sang phải, tới một điểm \(B\)
trên một đoạn nào đó thì phát hiện một đỉnh cao hơn ở bên trái, nên quay sang
đi về bên trái. Điểm rẽ \(B\) phải có những tính chất gì? Gọi đoạn chứa \(B\)
là \(P_1P_2\). Khi đó \(B\) thỏa mãn:
\begin{itemize}
  \item \(B\) là giao điểm giữa \(P_1P_2\) và đường thẳng nối hai điểm nào đó
  bên trái \(A\), ký hiệu là \(T_1,T_2\);
  \item các điểm trước \(A\) đều nằm phía dưới đường thẳng \(T_1T_2\), theo
  nghĩa không nghiêm ngặt;
  \item \(T_1,T_2\) là hai điểm kề nhau trên bao lồi trên của các điểm trước
  \(A\).
\end{itemize}

Điều này có nghĩa: một điểm quan sát có ý nghĩa \(B\) sẽ là giao điểm đầu
tiên giữa đường kéo dài của một cạnh nào đó trên bao lồi trên và đường gấp
khúc của dãy núi. Vì tổng số cạnh được sinh ra trong quá trình duy trì bao
lồi trên là \(O(n)\), số điểm quan sát có ý nghĩa cũng là
\(O(n)\).\footnote{Đây cũng là một tính chất gây phấn khích.}

Vậy làm thế nào nhanh chóng tìm ra các điểm quan sát đó?

Nếu đường thẳng \(T_1T_2\) cắt đoạn \(P_1P_2\), điều này có nghĩa là khi
\(P_2\) được thêm vào bao lồi, \(T_2\) sẽ bị bật ra; còn khi \(P_1\) được
thêm vào bao lồi thì \(T_2\) chưa bị bật ra. Vì thế trong lúc duy trì bao
lồi, ta chỉ cần lấy giao điểm giữa các đường thẳng liên quan tới điểm bị bật
ra và điểm vừa chèn, là có thể thu được toàn bộ điểm quan sát. Độ phức tạp
vẫn là \(O(n)\).

Bài toán với hữu hạn điểm quan sát đã được giải ở trên. Vì thế toàn bộ bài
toán được giải trọn vẹn trong thời gian \(O(n\log n)\).

\subsection{Tổng kết}

\subsubsection{Nguồn cảm hứng}

Cảm hứng của bài toán đến từ suy nghĩ về việc leo núi trong đời sống thực. Dùng
kiến thức thuật toán để phân tích một vấn đề thực tế không chỉ nâng cao năng
lực phân tích thuật toán, mà còn phản ánh mục đích cốt lõi của việc học là
``học để dùng''. Nó cũng có thể khơi dậy hứng thú đối với thi Tin học và làm
ta thấy vui trong quá trình đó. Hy vọng cách suy nghĩ này đem lại gợi ý cho
mọi người.

\subsubsection{Tư duy giải bài}

Nhìn lại quá trình giải bài: trước hết ta đơn giản hóa bài toán, biến vô hạn
thành hữu hạn; tiếp đó chia bài toán thành các phần nhỏ hơn và lần lượt giải
quyết. Đây đều là những phương pháp thường dùng trong quá trình làm bài.

Dưới đây là vài cảm nhận của tác giả về phương pháp giải bài, xin chia sẻ với
mọi người.

Trong quá trình giải bài thông thường, ta cần dựng một chiếc cầu giữa điều
kiện đã biết và vấn đề cần tìm. Có thể dùng công cụ đã học để suy luận xuôi
từ điều kiện đã biết, xem ta còn thu được gì. Cũng có thể suy luận ngược từ
vấn đề cần tìm, xem ta còn cần biết gì. Khi cây cầu suy luận xuôi và cây cầu
suy luận ngược có một điểm trùng nhau, cây cầu nối từ điều kiện tới bài toán
đã được dựng xong, và bài toán cũng được giải quyết.

Tuy nhiên phương pháp giải bài biến hóa muôn hình vạn trạng; không tồn tại bí
kíp vạn năng nào. Điều đáng mừng là: khi số bài bạn đã giải càng nhiều, khả
năng bạn giải được bài tiếp theo càng lớn. Vì vậy, chỉ cần chăm luyện tập và
chịu suy nghĩ, năng lực giải bài chắc chắn sẽ không ngừng bước lên những nấc
thang mới.

\section{Lấy số trên lưới}

\begin{center}
Wang Kangning\\
Trường THPT trực thuộc Đại học Sư phạm Đông Bắc, Cát Lâm
\end{center}

\subsection*{Tóm tắt}

Đây là báo cáo lời giải cho một bài toán thú vị trong đợt làm đề đối kháng
của đội tuyển quốc gia năm 2013.

\subsection{Đề bài}

\subsubsection{Mô tả bài toán}

Cho một bảng ô vuông \(N\times N\). Mỗi ô chứa một số; số ở hàng \(i\), cột
\(j\) bằng \(A_i\cdot B_j\).

Một người đi từ góc trên bên trái của một hình chữ nhật con tới góc dưới bên
phải của hình chữ nhật đó. Mỗi bước chỉ được đi sang phải hoặc đi xuống một
ô. Anh ta cộng các số đi qua trên đường đi và muốn tổng này nhỏ nhất.

Anh ta quyết định đi \(T\) lần. Lần thứ \(i\), anh đi từ ô hàng \(P_i\), cột
\(Q_i\) tới ô hàng \(R_i\), cột \(S_i\). Hãy hỏi tổng nhỏ nhất của các số đi
qua trong mỗi lần.

\subsubsection{Dữ liệu vào}

Dòng đầu tiên gồm hai số nguyên \(N,T\), lần lượt là kích thước bảng và số
lượng truy vấn.

Dòng thứ hai gồm \(N\) số nguyên \(A_1,A_2,\ldots,A_N\). Dòng này được sinh
ngẫu nhiên.

Dòng thứ ba gồm \(N\) số nguyên \(B_1,B_2,\ldots,B_N\). Dòng này cũng được
sinh ngẫu nhiên.

Tiếp theo là \(T\) dòng, mỗi dòng gồm bốn số nguyên
\(P_i,Q_i,R_i,S_i\), có ý nghĩa như trên và thỏa mãn
\(P_i\le R_i,\ Q_i\le S_i\).

\subsubsection{Dữ liệu ra}

In ra \(T\) dòng; mỗi dòng là một số nguyên, biểu diễn đáp án của truy vấn
tương ứng.

\subsubsection{Ví dụ vào}

\begin{verbatim}
3 4
1 2 3
2 3 4
1 1 1 1
1 3 1 3
1 1 3 3
1 1 3 1
\end{verbatim}

\subsubsection{Ví dụ ra}

\begin{verbatim}
2
4
29
12
\end{verbatim}

\subsubsection{Phạm vi dữ liệu}

\begin{itemize}
  \item Với \(20\%\) số test: \(N\le 100,\ T\le 100\).
  \item Với \(50\%\) số test: \(N\le 3000,\ T\le 1000\).
  \item Với \(70\%\) số test: \(N\le 50000,\ T\le 20000\).
  \item Với \(100\%\) số test:
  \(N\le 200000,\ T\le 50000,\ 1\le A_i,B_i\le 10^6\).
\end{itemize}

Bảo đảm mọi truy vấn đều hợp lệ. Ngoài ví dụ, tất cả các giá trị \(A_i,B_i\)
đều được sinh ngẫu nhiên.

\paragraph{Giới hạn tài nguyên.}
Giới hạn thời gian là \(2\) giây, giới hạn bộ nhớ là \(512\) MB.

\subsection{Hướng giải}

\subsubsection{Thuật toán một}

Với mỗi truy vấn, dùng quy hoạch động trực tiếp trên toàn bộ hình chữ nhật
của truy vấn.

Độ phức tạp thời gian là \(O(T\cdot N^2)\), độ phức tạp bộ nhớ là
\(O(N^2)\). Cách này chỉ kỳ vọng đạt khoảng \(20\) điểm.

\subsubsection{Thuật toán hai}

Thuật toán trên trả lời truy vấn quá chậm, nên ta cần tăng tốc phần truy vấn.

Xét tư tưởng chia để trị. Mỗi lần ta dùng đường giữa để chia hình chữ nhật
thành hai nửa, dùng quy hoạch động tính đáp án từ từng điểm trên đường giữa
tới từng điểm trong hình chữ nhật hiện tại, rồi đệ quy xử lý hai hình chữ
nhật ở hai bên đường giữa. Nếu điểm đầu và điểm cuối của một truy vấn nằm ở
hai phía của đường giữa, ta liệt kê điểm trên đường giữa mà đường đi đi qua;
nếu không, truy vấn được đẩy xuống một hình chữ nhật con.

Nếu lần nào cũng cắt theo đường thẳng đứng, đặt độ phức tạp tiền xử lý là
\(\operatorname{Time}(x)\), ta có
\[
  \operatorname{Time}(x)
  =2\operatorname{Time}(x/2)+O(N^2x).
\]
Khi đó độ phức tạp tiền xử lý là \(O(N^3\log N)\). Nếu luân phiên cắt theo
chiều ngang và chiều dọc, ta có
\[
  \operatorname{Time}(x)
  =2\operatorname{Time}(x/2)+O(x^3),
\]
nên độ phức tạp tiền xử lý là \(O(N^3)\).

Tổng độ phức tạp thời gian là \(O(N^3+TN)\), độ phức tạp bộ nhớ là
\(O(N^3)\). Chi phí tiền xử lý vẫn quá cao, vì vậy cần tìm phương án khác.

\subsubsection{Thuật toán ba}

Xét thuật toán chia khối. Nếu chia hình vuông thành các khối \(S\times S\),
ta tiền xử lý đáp án từ mỗi điểm tới mỗi điểm trên biên khối, đồng thời tiền
xử lý đáp án giữa mọi cặp điểm trong cùng một khối.

Khi trả lời truy vấn, nếu hai điểm không nằm trong cùng một khối, ta liệt kê
vị trí đầu tiên mà đường đi đi qua biên khối; nếu chúng nằm trong cùng một
khối thì trả lời trực tiếp.

Độ phức tạp thời gian là
\[
  O\left(N^2\left(S+\frac{N}{S^2}\right)
    +T\cdot\frac{N}{S}\right),
\]
và độ phức tạp bộ nhớ là
\[
  O\left(N^2\left(S+\frac{N}{S^2}\right)\right).
\]
Độ phức tạp này cũng còn quá cao.

\subsubsection{Cải tiến}

Một nguyên nhân chính khiến các thuật toán trên kém hiệu quả là chúng chưa
dùng tới điều kiện quan trọng trong đề: hai dãy \(A\) và \(B\) đều được sinh
ngẫu nhiên.

Qua thử nghiệm và quan sát, ta thấy số điểm rẽ của đường đi tối ưu không
nhiều. Hãy xét một điểm rẽ và hai đoạn đường đi kề với nó. Nếu điểm rẽ nằm ở
một cột có trọng số không nhỏ hơn một cột nào đó bên trái trong vùng đang
xét, đồng thời cũng không nhỏ hơn một cột nào đó bên phải, thì luôn tồn tại
một cách biến đổi đường đi không tệ hơn và làm mất điểm rẽ đó.

Cụ thể, giả sử trong hình minh họa của bài gốc, cột thứ năm trong khung có
trọng số không nhỏ hơn trọng số của cột ngoài cùng bên trái và cột ngoài cùng
bên phải của khung. Ta xét hai cách biến đổi đường đi quanh điểm rẽ:
\begin{itemize}
  \item Nếu trọng số của hàng phía trên trong vùng nhỏ hơn trọng số của hàng
  phía dưới, dùng biến đổi thứ nhất sẽ cho đáp án tốt hơn.
  \item Nếu trọng số của hàng phía trên lớn hơn trọng số của hàng phía dưới,
  dùng biến đổi thứ hai sẽ cho đáp án tốt hơn.
  \item Nếu hai trọng số hàng bằng nhau, dùng cách biến đổi nào thì đáp án
  cũng không xấu đi.
\end{itemize}

Như vậy ta có thể loại bỏ điểm rẽ đang xét.

Từ lập luận trên suy ra: với một điểm rẽ bất kỳ, trọng số của hàng hoặc cột
chứa nó phải nhỏ hơn tất cả trọng số của các hàng hoặc cột ở một phía trong
vùng truy vấn. Nói cách khác, nếu nhìn theo một phía của đoạn, điểm rẽ phải
là một kỷ lục nhỏ mới.

Với dữ liệu ngẫu nhiên, xác suất để giá trị thứ \(i\) nhỏ hơn mọi giá trị
đứng trước nó là \(1/i\). Vì vậy số lượng điểm rẽ kỳ vọng là
\[
  1+\frac12+\cdots+\frac1N=O(\log N).
\]

\subsubsection{Thuật toán bốn}

Ta tìm vét cạn mọi hàng và cột có thể chứa điểm rẽ. Sau đó dùng quy hoạch
động để giải truy vấn, trong đó chỉ những hàng và cột có khả năng chứa điểm
rẽ mới được xem là trạng thái hợp lệ.

Theo phân tích kỳ vọng ở trên, số hàng và số cột cần xét đều là
\(O(\log N)\); kỳ vọng của tích cũng bằng tích của các kỳ vọng, nên số trạng
thái kỳ vọng là \(O(\log^2 N)\). Ngoài ra, ta lưu tổng tiền tố của hai dãy
\(A\) và \(B\) để tăng tốc các phép chuyển trạng thái dọc theo một đoạn thẳng.

Độ phức tạp thời gian kỳ vọng là
\[
  O\left(T\cdot (N+\log^2 N)\right)=O(TN).
\]
Cách này kỳ vọng đạt khoảng \(50\) điểm.

\subsubsection{Thuật toán năm}

Nút thắt của thuật toán trước là việc tìm các hàng và cột có khả năng chứa
điểm rẽ. Bài toán con này có thể phát biểu như sau: trong một khoảng của một
dãy, tìm các số nhỏ hơn mọi số ở phía trước hoặc nhỏ hơn mọi số ở phía sau.
Chỉ cần biết, với mỗi vị trí, phần tử đầu tiên ở phía sau hoặc phía trước nhỏ
hơn nó là đủ. Việc này có thể thực hiện bằng hàng đợi đơn điệu.

Chẳng hạn, để tìm phần tử đầu tiên ở phía sau nhỏ hơn mỗi phần tử, ta quét
dãy từ phải sang trái. Với phần tử hiện tại, loại khỏi đuôi hàng đợi đơn điệu
mọi phần tử không nhỏ hơn nó. Khi đó phần tử đầu tiên ở phía sau nhỏ hơn nó
chính là phần tử ở đuôi hàng đợi hiện tại; sau đó đưa phần tử hiện tại vào
hàng đợi.

Tiền xử lý chỉ cần thời gian tuyến tính. Mỗi truy vấn sau đó được trả lời
trong thời gian kỳ vọng \(O(\log^2 N)\).

Độ phức tạp thời gian kỳ vọng là
\[
  O(N+T\log^2 N).
\]
Cách này kỳ vọng đạt \(100\) điểm. Tới đây bài toán đã được giải quyết trọn
vẹn.

\subsection{Lời cảm ơn}

Xin cảm ơn CCF đã cung cấp nền tảng học tập và trao đổi.

Xin cảm ơn cô Kong Weiling và thầy Gu Fangming vì sự chỉ dẫn.

Xin cảm ơn các bạn học đã giúp đỡ tôi.

\section{Thảo luận bài \emph{Two strings}}

\begin{center}
Luo Gan\\
Trường Học Quân Hàng Châu, Chiết Giang
\end{center}

\subsection{Mô tả bài toán}

Cho hai xâu \(A\) và \(B\). Có năm loại thao tác:
\begin{enumerate}
  \item Thêm một ký tự vào đầu xâu \(A\).
  \item Thêm một ký tự vào cuối xâu \(A\).
  \item Thêm một ký tự vào đầu xâu \(B\).
  \item Thêm một ký tự vào cuối xâu \(B\).
  \item Hỏi số lần xâu \(B\) hiện tại xuất hiện trong xâu \(A\) hiện tại.
\end{enumerate}

Bảo đảm mọi ký tự đều là chữ cái thường, và ban đầu hai xâu \(A,B\) đều không
rỗng.

\subsubsection{Dữ liệu vào}

Dòng thứ nhất và dòng thứ hai lần lượt là xâu \(A\) và xâu \(B\) ban đầu.
Dòng thứ ba chứa một số nguyên \(m\), biểu diễn số thao tác. Tiếp theo là
\(m\) dòng, mỗi dòng biểu diễn một thao tác. Số đầu tiên của mỗi dòng là một
số trong khoảng \(1\)--\(5\); nếu giá trị này khác \(5\), sau số đó sẽ có một
chữ cái thường.

\subsubsection{Dữ liệu ra}

Với mỗi truy vấn, in ra một số nguyên trên một dòng, biểu diễn số lần xâu
\(B\) xuất hiện trong xâu \(A\).

\subsubsection{Quy mô dữ liệu}

\begin{itemize}
  \item Với \(10\%\) dữ liệu, tổng độ dài cuối cùng của hai xâu \(A\) và
  \(B\) không vượt quá \(200\), số thao tác không vượt quá \(10\).
  \item Với \(30\%\) dữ liệu, tổng độ dài cuối cùng của hai xâu \(A\) và
  \(B\) không vượt quá \(2000\), số thao tác không vượt quá \(1000\).
  \item Với \(100\%\) dữ liệu, tổng độ dài cuối cùng của hai xâu \(A\) và
  \(B\) không vượt quá \(200000\), số thao tác không vượt quá \(200000\).
\end{itemize}

\subsection{Ví dụ}

\subsubsection{Dữ liệu vào}

\begin{verbatim}
ababc
a
7
5
4 b
5
3 a
1 a
5
5
\end{verbatim}

\subsubsection{Dữ liệu ra}

\begin{verbatim}
2
2
1
1
\end{verbatim}

\subsubsection{Giải thích ví dụ}

\begin{center}
\begin{tabular}{c c c c}
\textbf{Lần hỏi} & \textbf{Xâu \(A\)} & \textbf{Xâu \(B\)}
& \textbf{Số lần \(B\) xuất hiện trong \(A\)} \\
\hline
Lần \(1\) & \texttt{ababc} & \texttt{a} & \(2\) \\
Lần \(2\) & \texttt{ababc} & \texttt{ab} & \(2\) \\
Lần \(3\) & \texttt{aababc} & \texttt{aab} & \(1\) \\
Lần \(4\) & \texttt{aababc} & \texttt{aab} & \(1\)
\end{tabular}
\end{center}

\subsection{Thuật toán đơn giản}

\subsubsection{Mô phỏng}

Mô phỏng từng thao tác. Với mỗi truy vấn, trực tiếp liệt kê các vị trí mà
xâu \(B\) có thể xuất hiện trong xâu \(A\), rồi so sánh trực tiếp.

\subsubsection{Mô phỏng tối ưu bằng KMP}

Trên cơ sở mô phỏng, với mỗi truy vấn ta dùng thuật toán KMP để thực hiện
phép khớp xâu.

\subsection{Tối ưu mô phỏng bằng xử lý ngoại tuyến}

Đọc trước toàn bộ thao tác liên quan tới xâu \(A\), rồi ghép chúng thành một
``xâu \(A\) lớn''.

Khi đó, ở mỗi thời điểm truy vấn, xâu \(A\) hiện tại tương đương với một đoạn
nào đó của xâu \(A\) lớn. Bài toán trở thành hỏi trong đoạn đó có bao nhiêu
xâu \(B\).

Mỗi khi gặp thao tác loại \(3\) hoặc loại \(4\), ta chạy lại KMP một lần và
ghi nhận những vị trí khớp với xâu \(B\). Như vậy truy vấn được chuyển thành:
trong các vị trí đó, có bao nhiêu vị trí nằm trong đoạn từ \(L\) tới
\(R-\operatorname{Len}(B)+1\).

Vì \(L\) giảm dần còn \(R\) tăng dần, mỗi truy vấn có thể được giải trong
thời gian \(O(1)\).

\subsubsection{Độ phức tạp thời gian của từng thao tác}

\begin{center}
\begin{tabular}{c c}
\textbf{Thao tác} & \textbf{Độ phức tạp} \\
\hline
Thao tác \(1\) & \(O(n)\to O(1)\) \\
Thao tác \(2\) & \(O(n)\to O(1)\) \\
Thao tác \(3\) & \(O(n)\) \\
Thao tác \(4\) & \(O(n)\) \\
Thao tác \(5\) & \(O(n)\to O(1)\)
\end{tabular}
\end{center}

\subsection{Thuật toán xử lý ngoại tuyến bằng mảng hậu tố}

\subsubsection{Từ trực tuyến sang ngoại tuyến}

Sử dụng tư tưởng ngoại tuyến ở trên: lấy kết quả cuối cùng của hai xâu
\(A,B\), ghép chúng lại với nhau, đồng thời chèn một ký tự phân cách ở giữa,
chẳng hạn \texttt{+}. Sau đó dùng mảng hậu tố để duy trì quan hệ giữa các
hậu tố.

Ví dụ:

\begin{center}
\begin{tabular}{l}
Xâu \(A\) lớn: \texttt{ABABC} \\
Xâu \(B\) lớn: \texttt{ABA} \\
Xâu cuối cùng: \texttt{ABABC+ABA}
\end{tabular}
\end{center}

Sau khi sắp xếp mảng hậu tố:

\begin{center}
\begin{tabular}{l}
\texttt{A} \\
\texttt{ABA} \\
\texttt{ABC+ABA} \\
\texttt{ABABC+ABA} \\
\texttt{BA} \\
\texttt{BABC+ABA} \\
\texttt{BC+ABA} \\
\texttt{C+ABA} \\
\texttt{+ABA}
\end{tabular}
\end{center}

\subsubsection{Tìm khoảng hậu tố khả thi}

Xâu \(B\) hiện tại xuất hiện trong xâu \(A\) ở các vị trí từ \(i\) tới
\(i+\operatorname{Len}(B)-1\) khi và chỉ khi độ dài tiền tố chung dài nhất
\(Lcp\) của xâu \(B\) hiện tại và hậu tố thứ \(i\) của xâu \(A\) thỏa mãn
\[
  Lcp\ge \operatorname{Len}(B).
\]

Ta có thể dùng phương pháp nhân đôi để tiền xử lý mảng \texttt{height}, rồi
tìm khoảng hậu tố này trong thời gian \(O(\log n)\).

\subsubsection{Thống kê vị trí bắt đầu hợp lệ trong xâu A}

Ta định nghĩa một vị trí hợp lệ như sau: với vị trí \(i\) trong xâu \(A\)
lớn, toàn bộ đoạn
\[
  [i,\, i+\operatorname{Len}(B)-1]
\]
đều nằm trong xâu \(A\) hiện tại.

\subsubsection{Duy trì bằng cây BIT}

Ta có thể nhận thấy: mỗi thao tác bất kỳ chỉ làm thay đổi tính hợp lệ của
nhiều nhất hai vị trí. Vì vậy có thể dùng cây BIT để duy trì, trên một đoạn
của mảng hậu tố, có bao nhiêu vị trí hợp lệ.

\subsubsection{Độ phức tạp thời gian của từng thao tác}

\begin{center}
\begin{tabular}{c c}
\textbf{Thao tác} & \textbf{Độ phức tạp} \\
\hline
Thao tác \(1\) & \(O(\log n)\) \\
Thao tác \(2\) & \(O(\log n)\) \\
Thao tác \(3\) & \(O(\log n)\) \\
Thao tác \(4\) & \(O(\log n)\) \\
Thao tác \(5\) & \(O(\log n)\)
\end{tabular}
\end{center}

\subsection{Các điểm kiểm tra}

\begin{itemize}
  \item Suy nghĩ sâu về bản chất của bài toán.
  \item Nắm vững nhiều loại cấu trúc dữ liệu.
  \item Khả năng tối ưu nhiều loại thuật toán mô phỏng.
  \item Khả năng vận dụng thuật toán ngoại tuyến.
\end{itemize}

\section{\emph{Catch The Penguins}}

\begin{center}
Zhang Wentao\\
Trường Học Quân Hàng Châu, Chiết Giang
\end{center}

\subsection{Tóm tắt bài toán}

Cho tọa độ của \(N\) con chim cánh cụt trong không gian bốn chiều.

Có \(Q\) truy vấn. Mỗi truy vấn cho một điểm và hỏi có bao nhiêu chim cánh
cụt có tọa độ ở mọi chiều đều nhỏ hơn hoặc bằng tọa độ tương ứng của điểm
đó.

\paragraph{Ví dụ vào.}
\begin{verbatim}
1
0 0 0 0
2
1 1 1 0
1 1 1 -1
\end{verbatim}

\paragraph{Ví dụ ra.}
\begin{verbatim}
1
0
\end{verbatim}

\paragraph{Quy mô dữ liệu.}
Có tất cả \(20\) bộ dữ liệu.

\begin{center}
\begin{tabular}{c l}
\textbf{Dữ liệu} & \textbf{Ràng buộc} \\
\hline
1--2 & \(N,Q\le 5000\) \\
3--6 & \(N,Q\le 10000\) \\
7--14 & \(N\le 30000,\ Q\le 10000\) \\
15--18 & \(N\le 10000,\ Q\le 30000\) \\
19--20 & \(N,Q\le 30000\)
\end{tabular}
\end{center}

\subsection{Thuật toán 0}

Trực tiếp vét cạn. Với mỗi truy vấn, dùng một vòng lặp \texttt{for} để liệt
kê từng chim cánh cụt, rồi kiểm tra xem cả bốn tọa độ của nó có đều nhỏ hơn
hoặc bằng tọa độ của truy vấn hay không.

Độ phức tạp thời gian là \(O(n^2)\).

Kỳ vọng đạt \(10\) điểm.

\subsection{Tối ưu thuật toán 0}

Sắp xếp riêng các chim cánh cụt theo từng chiều trong bốn chiều tọa độ.

Với mỗi truy vấn, trên mỗi chiều ta tìm nhị phân trong mảng đã sắp xếp tương
ứng để biết số chim cánh cụt có tọa độ ở chiều đó nhỏ hơn hoặc bằng tọa độ
của truy vấn. Lấy giá trị nhỏ nhất trong bốn kết quả này để chọn chiều đem
ra liệt kê, như một tối ưu hằng số.

Kỳ vọng đạt \(30\) điểm.

\subsection{Tiếp tục tối ưu?}

Nếu muốn dùng thuật toán trực tuyến, dường như không có cấu trúc dữ liệu phù
hợp để trả lời truy vấn tọa độ bốn chiều, nên rất khó tiếp tục theo hướng đó.

Ta hãy thử chuyển sang xét cách xử lý truy vấn ngoại tuyến.

\subsection{Thuật toán chia khối}

Khi đã có ý tưởng ngoại tuyến, kết hợp thêm tư tưởng chia khối thì không khó
nghĩ ra một thuật toán có độ phức tạp tốt hơn một chút so với vét cạn.

\subsubsection{Rời rạc hóa và xử lý ngoại tuyến}

Sắp xếp riêng các truy vấn và các chim cánh cụt theo tọa độ chiều thứ nhất.

Sau đó liệt kê các truy vấn.

\subsubsection{Với truy vấn đang được liệt kê}

Thêm vào tất cả các chim cánh cụt có tọa độ chiều thứ nhất nhỏ hơn hoặc bằng
tọa độ chiều thứ nhất của truy vấn.

Sắp xếp theo chiều thứ hai. Mỗi khi thêm \(\sqrt n\) chim cánh cụt thì thực
hiện một lần sắp xếp bằng vét cạn; phần còn lại cũng xử lý vét cạn.

Độ phức tạp của phần này là \(O(n\sqrt n)\).

\subsubsection{Khi thực hiện một lần sắp xếp vét cạn}

Giả sử số chim cánh cụt đã được thêm vào và đã sắp xếp là \(M\).

Ta chia khối, sắp xếp theo tọa độ chiều thứ ba, rồi theo chiều thứ tư duy trì
cây đoạn lưu trọng số tiền tố.

Cụ thể, có thể dùng cây đoạn hàm thức.

\subsubsection{Truy vấn}

Với các khối đầy đủ: dùng nhị phân cộng truy vấn.

Với các phần không trọn khối: xử lý vét cạn.

Tổng độ phức tạp là \(O(n\sqrt n\log n)\).

Kỳ vọng đạt \(60\)--\(95\) điểm.

\subsection{Thuật toán tốt hơn}

Thông thường, xử lý chia khối là một phương án thay thế hiệu quả khi tồn tại
thuật toán có độ phức tạp tốt hơn nhưng độ phức tạp tư duy lại khá cao. Ở
đây, chỉ cần suy nghĩ thêm một chút là ta có thể tìm được một thuật toán chia
để trị tốt hơn.

\subsubsection{Bước 1}

Sắp xếp chung các truy vấn và chim cánh cụt theo tọa độ chiều thứ nhất.

Sau đó thực hiện ``trộn'' theo tọa độ chiều thứ hai.

\subsubsection{Bước 2}

Giả sử dãy chỉ số sau khi sắp xếp là \(\texttt{id[L]}\) tới
\(\texttt{id[R]}\), trong đó có phần tử là chim cánh cụt và có phần tử là
truy vấn. Gọi \(\texttt{MID}\) là điểm cuối của đoạn bên trái.

\subsubsection{Bước 3}

Đệ quy xử lý đáp án của hai phần \(\texttt{L}\)--\(\texttt{MID}\) và
\(\texttt{MID+1}\)--\(\texttt{R}\). Dĩ nhiên, khi \(\texttt{L}\) và
\(\texttt{R}\) bằng nhau thì có thể trả về ngay.

\subsubsection{Bước 4}

Ta thấy trường hợp chưa được xử lý chỉ còn là ảnh hưởng của chim cánh cụt ở
nửa trái, tức \(\texttt{L}\)--\(\texttt{MID}\), lên truy vấn ở nửa phải, tức
\(\texttt{MID+1}\)--\(\texttt{R}\). Lúc này cả hai nửa đều đã được sắp xếp
theo tọa độ chiều thứ hai (lý do được giải thích ngay sau đây), và thứ tự
theo chiều thứ nhất của chúng không còn cần thiết nữa.

Ta trộn nửa trái và nửa phải theo tọa độ chiều thứ hai, tức gộp hai bảng
tuyến tính đã có thứ tự. Khi tạo bảng mới, nếu phần tử được thêm là chim
cánh cụt ở nửa trái, ta chèn điểm tạo bởi tọa độ chiều thứ ba và chiều thứ tư
của nó vào một ``tập hợp''. Nếu phần tử được thêm là truy vấn ở nửa phải, ta
hỏi trong ``tập hợp'' có bao nhiêu điểm nằm ở góc dưới bên trái của điểm tạo
bởi hai tọa độ chiều thứ ba và chiều thứ tư của truy vấn đó, rồi cộng số này
vào đáp án của truy vấn.

\subsubsection{Bước 5}

Về cách cài đặt ``tập hợp'', ta có thể dùng cây BIT (hoặc cây đoạn) lồng cây
cân bằng.

Tổng độ phức tạp thời gian là \(O(n\log^3 n)\).

\subsection{Các điểm kiểm tra}

\begin{itemize}
  \item Chuyển đổi tư duy từ trực tuyến sang ngoại tuyến.
  \item Vận dụng các tư tưởng chia khối, chia để trị.
  \item Xử lý thứ tự dữ liệu một cách khéo léo.
  \item Nắm vững cấu trúc dữ liệu.
\end{itemize}

\section{Bàn về ứng dụng tư tưởng chia khối trong một lớp bài xử lý dữ liệu}

\begin{center}
Luo Jianqiao\\
Trường THPT số 80 Bắc Kinh
\end{center}

\subsection*{Tóm tắt}

Trong thi đấu và thực tiễn, làm thế nào để tổ chức và xử lý dữ liệu quy mô
lớn một cách hiệu quả là một vấn đề quan trọng. Bài viết này chủ yếu thảo
luận một lớp bài xử lý dữ liệu trong thi đấu, liên quan tới dãy tuyến tính và
cấu trúc cây. Tác giả đưa vào khái niệm tư tưởng chia khối, khảo sát ứng dụng
của nó trong tổ chức dữ liệu, tiền xử lý dữ liệu và thiết kế thuật toán ngoại
tuyến; đồng thời dùng ví dụ để minh họa cách chia dữ liệu thành khối nhằm thu
được thuật toán có độ phức tạp thời gian thấp.

\subsection{Dẫn nhập}

Trong một lớp bài xử lý dữ liệu, ta cần thường xuyên thống kê thông tin về
những phần tử thỏa mãn điều kiện nào đó trong một tập dữ liệu rất lớn, chẳng
hạn số lượng hoặc giá trị cực trị. Các cấu trúc dữ liệu truyền thống, như cây
đoạn, tổ chức toàn bộ dữ liệu trong tập theo một cách có thứ tự, để mỗi truy
vấn có thể được tách thành tổng thông tin của một số phần có thứ tự, qua đó
nâng cao hiệu quả.

Bài viết này giới thiệu một tư tưởng chia khối. Nó cũng làm dữ liệu có thứ tự
và có tầng bậc, nhưng phương thức khác đi, khả năng mở rộng mạnh hơn và phạm
vi áp dụng rộng hơn.

Cốt lõi của tư tưởng chia khối là chia một tập hợp thành một số tập con có
quy mô nhỏ hơn. Nói chung, ta có xu hướng đưa các phần tử có tính chất gần
nhau vào cùng một tập con, rồi tổ chức các phần tử trong mỗi tập con theo một
cấu trúc nhất định.

Trong các bài xử lý dữ liệu, tư tưởng chia khối có phạm vi ứng dụng rộng vì
cách chia khối có vài tính chất tốt:
\begin{itemize}
  \item Quy mô của mỗi tập con nhỏ, nên có thể dùng thuật toán có độ phức tạp
  cao theo quy mô tập con.
  \item Số lượng tập con ít, nên có thể dùng thuật toán liên quan tới quy mô
  toàn cục để xử lý riêng từng tập con.
  \item Các phần tử trong mỗi tập con có tính chất chung, nên có thể thiết kế
  thuật toán có tính nhắm đích cho từng tập con.
\end{itemize}

Dễ thấy hai tính chất đầu mâu thuẫn với nhau: để cân bằng, ta không thể để
tập con quá lớn, nhưng cũng không thể để số tập con quá lớn. Đây là lý do độ
phức tạp thời gian của các thuật toán chia khối thường liên quan tới
\(\sqrt N\).

Bài viết chủ yếu bàn về ứng dụng của chia khối trong các bài xử lý dữ liệu
liên quan tới dãy tuyến tính và cấu trúc cây. Mục 2 giới thiệu phương pháp
chia khối chung trong các cấu trúc dữ liệu điển hình. Mục 3 giới thiệu ứng
dụng của chia khối trong tiền xử lý dữ liệu. Mục 4 giới thiệu ứng dụng của
chia khối khi thiết kế thuật toán ngoại tuyến.

\subsection{Chia khối dữ liệu}

Mục này giới thiệu phương pháp chia khối chung cho hai cấu trúc dữ liệu
thường gặp là dãy tuyến tính và cây, rồi lần lượt đưa ra ví dụ dùng chia khối
để giải.

\subsubsection{Chia khối dãy tuyến tính}

Trong dãy tuyến tính, mọi phần tử của tập dữ liệu được sắp xếp theo một thứ
tự nhất định. Một bài điển hình là thường xuyên thống kê ``tổng'' thông tin
của các phần tử thỏa mãn một tính chất nào đó trong một dãy con liên tiếp; ở
đây ``tổng'' nghĩa là đáp án có thể được tách trực tiếp hoặc gián tiếp thành
tổng đáp án của các phần khác nhau. Trước hết ta giới thiệu một cách chia
khối thường dùng, sau đó minh họa cách sử dụng nó để giải hiệu quả loại bài
này.

Cách chia khối như sau: bắt đầu từ phần tử đầu tiên của dãy, cứ mỗi \(S\)
phần tử liên tiếp tạo thành một khối. Nếu cuối cùng còn ít hơn \(S\) phần tử,
cho các phần tử còn lại tạo thành một khối.

\paragraph{Tính chất 2.1.}
Giả sử \(N\) là độ dài dãy. Khi dùng cách trên, số khối không vượt quá
\(\lceil N/S\rceil+1\).

\paragraph{Tính chất 2.2.}
Giả sử một truy vấn hỏi thông tin của dãy con liên tiếp từ phần tử thứ \(L\)
tới phần tử thứ \(R\), gọi tắt là đoạn \([L,R]\). Khi đó, ngoài các phần tử
nằm trong những khối được chứa trọn vẹn, số phần tử của đoạn \([L,R]\) còn
dính tới không vượt quá \(2S\).

Theo hai tính chất trên, nếu ta có thể duy trì ở tầng khối tổng thông tin của
toàn bộ phần tử trong một khối, thì chỉ cần chọn kích thước khối \(S\) thích
hợp, với mọi đoạn ta đều có thể nhanh chóng tìm ra những khối trọn vẹn và
những phần tử thừa, rồi gộp chúng để thu được thông tin của cả đoạn.

\paragraph{Ví dụ một.}
Cho một dãy gồm \(N\) số nguyên. Cần thực hiện \(M\) thao tác thuộc hai loại:
\begin{itemize}
  \item \texttt{ADD D\_i X\_i}: bắt đầu từ số thứ nhất, cứ cách \(D_i\) vị
  trí thì cộng thêm \(X_i\) vào phần tử đó.
  \item \texttt{QUERY L\_i R\_i}: hỏi tổng các phần tử từ vị trí \(L_i\) tới
  \(R_i\).
\end{itemize}

Quy mô dữ liệu: \(1\le N,M\le 10^5\); ở mọi thời điểm, giá trị tuyệt đối của
mỗi số trong dãy không vượt quá \(10^9\).

\paragraph{Phân tích thuật toán.}
Trong trường hợp xấu nhất, mỗi thao tác \texttt{ADD} và \texttt{QUERY} đều
có thể liên quan tới \(N\) phần tử, nên mô phỏng đơn thuần sẽ rất chậm. Ta
xét chia dãy thành khối và ghi ở tầng khối tổng các phần tử trong từng khối.

\paragraph{Bổ đề 2.1.}
Nếu đặt kích thước khối là \(\sqrt N\), mỗi truy vấn có thể xử lý trong
\(O(\sqrt N)\) thời gian.

Cách làm cụ thể là khi truy vấn, liệt kê từng khối; nếu khối nằm trọn trong
đoạn truy vấn thì cộng trực tiếp tổng của khối vào đáp án. Số phần tử còn lại
không vượt quá \(2\sqrt N\), nên có thể liệt kê trực tiếp.

Tuy nhiên, độ phức tạp của thao tác sửa vẫn chưa được cải thiện. Lúc này ta
tiếp tục dùng tư tưởng chia khối và chú ý tới sự thật sau:

\paragraph{Bổ đề 2.2.}
Khi và chỉ khi thao tác \texttt{ADD} có \(D_i<\sqrt N\), số phần tử cần cập
nhật mới vượt quá \(\sqrt N\).

Vì vậy ta phân loại thao tác \texttt{ADD}. Nếu \(D_i\ge\sqrt N\), cập nhật
trực tiếp các phần tử bị ảnh hưởng và tổng khối tương ứng, độ phức tạp là
\(O(\sqrt N)\). Nếu \(D_i<\sqrt N\), chỉ cần dùng mảng để ghi tổng lượng sửa
của mọi thao tác \texttt{ADD} có cùng \(D_i\); lúc này một thao tác sửa là
\(O(1)\).

Khi đó, giá trị của từng phần tử và tổng từng khối chỉ xét loại sửa thứ nhất.
Khi truy vấn, ta còn phải liệt kê mọi \(D_i\) thuộc loại thứ hai và tính ảnh
hưởng của tổng lượng sửa với \(D_i\) đó lên đáp án. Vì số \(D_i\) thuộc loại
thứ hai không vượt quá \(\sqrt N\), độ phức tạp một truy vấn vẫn là
\(O(\sqrt N)\).

\paragraph{Định lý 2.3.}
Bài toán trên có thể giải trong thời gian \(O((N+M)\sqrt N)\).

\subsubsection{Chia khối cấu trúc cây}

Trong cấu trúc cây, các đỉnh kề nhau có tương quan rất mạnh. Một ý tưởng trực
quan là mong muốn chia cây thành một số khối liên thông có quy mô nhỏ. Dưới
đây là một cách chia khối trên cây: với \(S\) thích hợp, nó có thể chia cây
thành \(O(N/S)\) khối liên thông kích thước \(O(S)\), trong đó \(N\) là số
đỉnh của cây.

Trước hết cần khái niệm dãy DFS. Với cây vô căn, có thể chọn tùy ý một đỉnh
làm gốc để biến nó thành cây có gốc.

\paragraph{Định nghĩa 2.1 (Dãy DFS của cây).}
Duy trì một dãy, ban đầu rỗng. Từ gốc, thực hiện duyệt sâu; mỗi khi gặp một
đỉnh lúc vào stack hoặc ra khỏi stack thì thêm nó vào cuối dãy. Dãy thu được
sau khi duyệt xong gọi là dãy DFS.

Như vậy, với một cây bất kỳ ta thu được một dãy DFS có \(2N\) phần tử; mỗi
đỉnh của cây tương ứng với hai phần tử trong dãy. Dãy DFS có những tính chất
tốt.

\paragraph{Bổ đề 2.4.}
Trong dãy DFS, quan hệ giữa hai phần tử kề nhau chỉ có ba loại: hoặc là cùng
một đỉnh, hoặc là quan hệ cha con, hoặc là quan hệ anh em. Loại thứ ba xảy ra
khi và chỉ khi trong DFS, phần tử trước là đỉnh được đưa ra khỏi stack, còn
phần tử sau là đỉnh vừa được đưa vào stack.

Từ bổ đề 2.4 có thể chứng minh tính chất sau.

\paragraph{Bổ đề 2.5.}
Một dãy con liên tiếp từ hạng \(u\) tới hạng \(v\) trong dãy DFS vừa đúng
tương ứng với một đường đi từ đỉnh ứng với hạng \(u\) tới đỉnh ứng với hạng
\(v\). Ngoài LCA của hai đầu mút ra, mọi đỉnh trên đường đi đều xuất hiện ít
nhất một lần trong dãy con đó.

Từ đó có kết luận ta cần:

\paragraph{Định lý 2.6.}
Dùng phương pháp chia khối của mục trước để chia dãy DFS của cây thành các
khối kích thước \(S\). Với mỗi khối của dãy, tập các đỉnh tương ứng với toàn
bộ phần tử trong khối cùng LCA của chúng tạo thành một khối liên thông trên
cây. Số khối liên thông không vượt quá \(2N/S\).

Phương pháp chia khối này cài đặt đơn giản, đồng thời có tính chất chung của
tư tưởng chia khối: mỗi khối không lớn và số khối cũng không nhiều. Đáng tiếc
là có thể tồn tại một số đỉnh xuất hiện trong nhiều hơn hai khối liên thông.

\paragraph{Ví dụ hai.}
Cho một cây \(N\) đỉnh và một số nguyên \(K\). Mỗi cạnh có trọng số. Cần tính
với mỗi đỉnh \(i\), trong các khoảng cách từ \(N-1\) đỉnh còn lại tới đỉnh
\(i\), giá trị nhỏ thứ \(K\) là bao nhiêu.

Quy mô dữ liệu: \(1\le K<N\le 50000\); trọng số cạnh là số nguyên có giá trị
tuyệt đối nhỏ hơn \(1000\).

\paragraph{Phân tích thuật toán.}
Thuật toán đơn giản là lần lượt chọn mỗi đỉnh làm gốc, tính khoảng cách từ
mọi đỉnh khác tới nó, rồi sắp xếp để lấy đáp án. Độ phức tạp thời gian là
\(O(N^2\log N)\). Nếu dùng sắp xếp cơ số, có thể giảm xuống \(O(N^2)\).

Ta chia cây theo phương pháp ở trên thành một số khối liên thông kích thước
\(O(S)\). Xét việc lợi dụng tính liên quan giữa các đỉnh kề nhau để nhanh
chóng tính đáp án khi từng đỉnh trong cùng một khối liên thông được chọn làm
gốc.

\paragraph{Bổ đề 2.7.}
Giả sử \(u\) là một đỉnh bất kỳ không nằm trong khối liên thông hiện tại. Dù
chọn đỉnh nào trong khối làm gốc, trên đường từ \(u\) tới gốc, đỉnh đầu tiên
nằm trong khối là không đổi.

Theo bổ đề 2.7, ta phân loại mọi đỉnh ngoài khối theo đỉnh đầu tiên trong
khối mà đường đi tới gốc đi qua, gọi tắt là tổ tiên gần nhất trong khối. Rõ
ràng các đỉnh ngoài khối được chia thành nhiều nhất \(S\) loại.

\paragraph{Bổ đề 2.8.}
Giả sử \(u\) và \(v\) là hai đỉnh bất kỳ không nằm trong khối hiện tại và
thuộc cùng một loại. Gọi \(p\) là tổ tiên gần nhất trong khối của chúng. Khi
chọn bất kỳ đỉnh \(r\) nào trong khối làm gốc, quan hệ thứ tự giữa khoảng
cách từ \(u\) tới \(r\) và từ \(v\) tới \(r\) đúng bằng quan hệ thứ tự giữa
khoảng cách từ \(u\) tới \(p\) và từ \(v\) tới \(p\).

Vì vậy, với mỗi loại, ta sắp xếp các đỉnh của loại đó theo khoảng cách tới
tổ tiên gần nhất trong khối từ nhỏ tới lớn. Sau đó liệt kê từng đỉnh trong
khối làm gốc và lần lượt tính đáp án.

\paragraph{Định lý 2.9.}
Khi chọn một đỉnh \(r\) bất kỳ trong khối liên thông làm gốc, có thể trong
\(O(S\log^2 N)\) thời gian tính được giá trị nhỏ thứ \(K\) trong các khoảng
cách từ các đỉnh khác tới \(r\).

Cách làm cụ thể như sau. Để tính giá trị nhỏ thứ \(K\), ta nhị phân đáp án.
Giả sử giá trị nhị phân là \(X\). Từ \(r\), liệt kê mọi đỉnh trong khối, tính
khoảng cách của chúng tới \(r\) và đếm có bao nhiêu khoảng cách không vượt
quá \(X\). Với mỗi đỉnh \(p\) trong khối, giả sử khoảng cách từ \(p\) tới
\(r\) là \(D_p\); trong bảng có thứ tự của loại nhận \(p\) làm tổ tiên gần
nhất trong khối, nhị phân vị trí lớn nhất có giá trị không vượt quá
\(X-D_p\), qua đó tính số đỉnh ngoài khối thuộc loại này và cách \(r\) không
quá \(X\).

Sau khi xử lý mọi \(p\), giả sử số đỉnh cách \(r\) không quá \(X\) là
\(\texttt{cnt}\). Nếu \(\texttt{cnt}<K\), đáp án thật lớn hơn \(X\); ngược
lại đáp án thật không lớn hơn \(X\). Mỗi lần nhị phân thu hẹp một nửa phạm
vi đáp án, do đó đạt độ phức tạp \(O(S\log^2 N)\).

\paragraph{Bổ đề 2.10.}
Tính cả thời gian sắp xếp và số đỉnh trong khối, ta có thể tính đáp án cho
mọi đỉnh trong một khối liên thông trong
\[
  O(N\log N+S^2\log^2 N)
\]
thời gian.

\paragraph{Định lý 2.11.}
Nếu chọn kích thước khối \(S=\Theta(\sqrt{N/\log N})\), độ phức tạp thời gian
tổng cộng của thuật toán là \(O(N^{1.5}\log^{1.5}N)\).

\subsubsection{Mở rộng thêm}

Bạn đọc quan tâm có thể tự suy nghĩ cách mở rộng phương pháp chia khối dãy
tuyến tính lên trường hợp \(d\) chiều, cũng như việc phương pháp chia khối
cây có áp dụng được cho các cấu trúc phức tạp hơn hay không.

\subsection{Chia khối và tiền xử lý}

Tiền xử lý nghĩa là khi duy trì một tập dữ liệu, trước mọi thao tác truy vấn
ta xử lý dữ liệu trước, thường ghi lại nhiều thông tin phụ, rồi khi truy vấn
thì dùng thông tin đã xử lý để tăng hiệu quả thời gian. Nếu có thao tác sửa,
còn phải xét ảnh hưởng của sửa lên thông tin đã tiền xử lý.

Kết hợp tư tưởng chia khối với tiền xử lý có thể giải hiệu quả một lớp bài
truy vấn tập hợp rất rộng. Mục này trước hết nêu các tính chất mà loại bài có
thể tối ưu bằng tiền xử lý cần thỏa mãn, sau đó giới thiệu chi tiết phương
pháp tiền xử lý chung trên dãy và trên cây.

\subsubsection{Điều kiện để tiền xử lý}

Giả sử tập dữ liệu đang duy trì là tập \(S\).

\paragraph{Điều kiện 3.1 (Có thể tăng lượng).}
Giả sử hai lần truy vấn tương ứng với hai tập con \(A\) và \(B\) của tập
\(S\). Độ phức tạp để chuyển thông tin liên quan của tập truy vấn \(A\) thành
thông tin liên quan của tập truy vấn \(B\) là một hàm theo kích thước hiệu
đối xứng của \(A\) và \(B\), tức
\[
  (A-B)\cup(B-A).
\]
Nói cách khác, độ phức tạp chuyển trạng thái chỉ phụ thuộc vào số phần tử
thuộc hiệu đối xứng.

Gọi \(\mathcal{P}\) là tập hợp mọi tập con của \(S\) có thể xuất hiện trong
truy vấn.

\paragraph{Điều kiện 3.2 (Có thể tiền xử lý).}
Với các số nguyên \(K\) và \(L\) cho trước, tồn tại một tập con
\(\mathcal{Q}\subseteq \mathcal{P}\) có kích thước không vượt quá \(K\), sao
cho với mọi phần tử \(X\in\mathcal{P}\), tồn tại một phần tử
\(Y\in\mathcal{Q}\) mà kích thước hiệu đối xứng giữa \(X\) và \(Y\) không
vượt quá \(L\).

Nếu truy vấn trên tập dữ liệu thỏa mãn điều kiện 3.1, đồng thời tồn tại
\(K,L\) thích hợp để tập dữ liệu thỏa mãn điều kiện 3.2, ta có thể dùng
phương pháp chia khối và tiền xử lý để tăng hiệu quả thời gian của thao tác
truy vấn.

Khi đó ta tìm tập \(\mathcal{Q}\) trong điều kiện 3.2 và tiền xử lý thông tin
liên quan của các phần tử trong \(\mathcal{Q}\). Mỗi khi truy vấn phần tử
\(X\), tìm phần tử \(Y\in\mathcal{Q}\) tương ứng. Vì thông tin của \(Y\) đã
được tiền xử lý, và truy vấn thỏa mãn điều kiện 3.1, ta có thể nhanh chóng
chuyển thông tin của \(Y\) thành thông tin của \(X\) trong
\(O(f(L))\) thời gian.

\subsubsection{Tiền xử lý trên dãy tuyến tính}

Mục nhỏ này xét một lớp bài trên dãy tuyến tính: thường xuyên truy vấn thông
tin của một dãy con liên tiếp và thỏa mãn điều kiện tiền xử lý. Trước hết là
phương pháp tiền xử lý, sau đó là một ví dụ.

Ta chia dãy thành khối. Giả sử kích thước mỗi khối là \(S\), số khối là
\(O(N/S)\). Theo điều kiện 3.1, với bất kỳ đoạn có \(\texttt{len}\) phần tử,
ta có thể thu được thông tin liên quan trong thời gian
\(O(f(\texttt{len}))\). Do đó ta tiền xử lý thông tin liên quan giữa hai khối
bất kỳ như sau.

Lấy một khối \(A\) bất kỳ, xét việc lần lượt tính thông tin giữa khối \(A\)
và từng khối \(B\) sau nó theo thứ tự vị trí trong dãy. Sau khi tính được
thông tin từ khối \(A\) tới khối \(B_i\), có thể trong \(O(f(S))\) thời gian
chuyển thành thông tin từ khối \(A\) tới khối \(B_{i+1}\).

\paragraph{Định lý 3.1.}
Chỉ xét các khối \(B_i\) nằm sau khối \(A\), nếu không tính độ phức tạp ghi
thông tin, có thể trong \(O(Nf(S)/S)\) thời gian tính thông tin giữa khối
\(A\) và mọi khối \(B_i\).

Nếu độ phức tạp ghi thông tin giữa hai khối bất kỳ có thể bỏ qua, ta liệt kê
khối \(A\) và hoàn thành toàn bộ quá trình chia khối, tiền xử lý trong
\(O(N^2f(S)/S^2)\) thời gian.

Theo tính chất 2.2, thông tin của một đoạn bất kỳ \([L,R]\) có thể được suy
ra từ thông tin của một dãy con giữa hai khối đã tiền xử lý bằng cách thêm
không quá \(2S\) phần tử. Vì vậy:

\paragraph{Định lý 3.2.}
Với bài thống kê đoạn trên dãy tuyến tính, nếu bài toán thỏa mãn điều kiện
tiền xử lý, trước hết chia khối và tiền xử lý dãy; sau đó mỗi truy vấn đoạn
\([L,R]\) có thể trả lời trong \(O(f(S))\) thời gian.

\paragraph{Ví dụ ba.}
Cho một dãy gồm \(N\) số, mỗi số là một số nguyên trong khoảng \(1\) tới
\(N\). Cần trả lời \(M\) truy vấn. Mỗi truy vấn cho ba số nguyên
\(L_i,R_i,K_i\), hỏi trong đoạn \([L_i,R_i]\) có bao nhiêu số xuất hiện đúng
\(K_i\) lần.

Quy mô dữ liệu: \(1\le N,M\le 10^5\).

\paragraph{Phân tích thuật toán.}
Trong bài này, đáp án của các đoạn kề nhau không thể nhanh chóng gộp thành
đáp án của đoạn hợp bằng cách ghi vài thông tin phụ, nên không thể dùng những
cấu trúc dữ liệu như cây đoạn để duy trì. Nhưng truy vấn có tính chất sau:

\paragraph{Bổ đề 3.3.}
Với hai đoạn bất kỳ \([L_1,R_1]\) và \([L_2,R_2]\), nếu ta đã ghi số lần xuất
hiện trong \([L_1,R_1]\) của từng số từ \(1\) tới \(N\), cũng như số lượng
các số có tần suất \(1\) tới \(N\), thì có thể trong
\[
  O(|L_1-L_2|+|R_1-R_2|)
\]
thời gian chuyển chúng thành thông tin tương ứng của đoạn \([L_2,R_2]\).

Vì vậy có thể dùng phương pháp tiền xử lý. Chia dãy thành khối với kích thước
mỗi khối là \(O(\sqrt N)\), số khối cũng là \(O(\sqrt N)\). Ta có thể trong
\(O(N\sqrt N)\) thời gian tính số lần xuất hiện của mỗi số từ \(1\) tới
\(N\), và số lượng các số có tần suất \(1\) tới \(N\), giữa hai khối bất kỳ.
Nhưng nếu ghi toàn bộ thông tin như vậy thì dung lượng là \(O(N^2)\), nên cần
giảm lượng thông tin phải ghi.

\paragraph{Bổ đề 3.4.}
Đánh số các khối từ trái sang phải là \(1\) tới \(\sqrt N\). Số lần xuất hiện
của một số \(x\) trong các khối từ \(i\) tới \(j\) bằng số lần \(x\) xuất
hiện trong \(j\) khối đầu trừ số lần \(x\) xuất hiện trong \(i-1\) khối đầu.

Ta chỉ duy trì số lần xuất hiện của mỗi số từ \(1\) tới \(N\) trong tiền tố
từ \(1\) tới \(\sqrt N\) khối, giảm độ phức tạp không gian xuống
\(O(N\sqrt N)\). Việc này tương đương với duy trì số lần xuất hiện giữa hai
khối bất kỳ.

\paragraph{Bổ đề 3.5.}
Số giá trị xuất hiện trong dãy ít nhất \(\sqrt N\) lần không vượt quá
\(\sqrt N\).

Khi tiền xử lý, ta chỉ ghi số lượng các giá trị có số lần xuất hiện từ \(1\)
tới \(\sqrt N\) giữa hai khối. Do đó độ phức tạp ghi thông tin là
\(O(N\sqrt N)\). Theo định lý 3.2, ta có thể trả lời đúng mọi truy vấn có
\(K_i<\sqrt N\) trong \(O(\sqrt N)\) thời gian.

Với truy vấn có \(K_i\ge\sqrt N\), ta xử lý trước tổng tiền tố tần suất của
mọi giá trị xuất hiện nhiều hơn \(\sqrt N\) lần. Khi truy vấn, liệt kê những
giá trị này trong \(O(\sqrt N)\) thời gian, và \(O(1)\) để kiểm tra tần suất
của từng giá trị trong đoạn truy vấn có bằng \(K_i\) hay không.

\paragraph{Định lý 3.6.}
Thuật toán trên có thể hoàn thành tiền xử lý trong \(O(N\sqrt N)\) thời gian.
Mỗi truy vấn được trả lời trong \(O(\sqrt N)\) thời gian. Tổng độ phức tạp
thời gian và không gian là \(O((N+M)\sqrt N)\).

\subsubsection{Tiền xử lý trên cấu trúc cây}

Mục nhỏ này xét một lớp bài trên cây: thường xuyên truy vấn thông tin của
đường đi giữa hai đỉnh và thỏa mãn điều kiện tiền xử lý. Ta mở rộng phương
pháp chia khối tiền xử lý từ dãy tuyến tính lên cây, mô tả cách đánh dấu các
đỉnh then chốt trên cây, rồi đưa ra một ví dụ.

\paragraph{Bổ đề 3.7.}
Trong cây, đường đi đơn giữa hai đỉnh bất kỳ là duy nhất.

\paragraph{Bổ đề 3.8.}
Với hai đỉnh bất kỳ \(u,v\) trong cây có gốc, gọi \(p\) là LCA của chúng.
Đường đi từ \(u\) tới \(v\) vừa đúng bằng đường đi từ \(u\) tới \(p\) và
đường đi từ \(p\) tới \(v\).

Ta muốn đánh dấu một số đỉnh của cây làm đỉnh then chốt, đồng thời tiền xử
lý đường đi giữa mọi cặp đỉnh then chốt, sao cho mọi đường đi đủ dài trên
cây đều chứa một đường đi đã được xử lý, còn phần còn lại ngoài đường đi đó
có số đỉnh ít. Dưới đây là một thuật toán tham lam.

Với cây có gốc, liệt kê các đỉnh theo thứ tự độ sâu giảm dần và quyết định
có đánh dấu chúng làm đỉnh then chốt hay không. Đánh dấu một đỉnh \(u\) khi
và chỉ khi trong cây con của \(u\) tồn tại một đỉnh mà trên đường từ đỉnh đó
tới \(u\) chưa có đỉnh then chốt, đồng thời số đỉnh trên đường này lớn hơn
\(S\). Ở đây \(S\) là số cho trước.

\paragraph{Bổ đề 3.9.}
Giả sử \(N\) là số đỉnh của cây. Thuật toán tham lam đánh dấu không quá
\(N/S\) đỉnh then chốt.

Lý do là mỗi đỉnh then chốt đều là đỉnh then chốt đầu tiên trên đường tới gốc
của ít nhất \(S\) đỉnh không then chốt.

\paragraph{Bổ đề 3.10.}
Giả sử độ sâu của gốc là \(0\). Với bất kỳ đỉnh nào có độ sâu không nhỏ hơn
\(S\), trong \(S+1\) tổ tiên gần nhất của nó, kể cả chính nó, có ít nhất một
đỉnh then chốt.

\paragraph{Hệ quả 3.11.}
Với bất kỳ đường đi đơn nào trong cây có số đỉnh lớn hơn \(3S\), gọi \(u\)
và \(v\) là hai đầu mút. Khoảng cách từ \(u\) tới đỉnh then chốt gần nhất
trên đường đi không vượt quá \(2S-2\); tổng khoảng cách từ \(u\) và \(v\) tới
các đỉnh then chốt gần nhất trên đường đi không vượt quá \(3S-3\).

Trong thuật toán tham lam đánh dấu đỉnh then chốt, chỉ cần ghi với mỗi đỉnh
giá trị lớn nhất của số đỉnh trên một đường trong cây con của nó, chỉ gồm các
đỉnh không then chốt và kết thúc tại nó, là có thể quyết định có đánh dấu nó
hay không.

\paragraph{Định lý 3.12.}
Trong \(O(N)\) thời gian, có thể đánh dấu \(O(N/S)\) đỉnh then chốt trong
cây. Nếu việc tiền xử lý giữa các cặp đỉnh then chốt có độ phức tạp
\(O((N/S)^2f(N))\), thì với mọi đường đi đơn, sau khi bỏ phần đã được xử lý,
số đỉnh còn lại là \(O(S)\).

\paragraph{Ví dụ bốn.}
Cho một cây \(N\) đỉnh. Màu của mỗi đỉnh là một số nguyên từ \(1\) tới \(M\).

Định nghĩa lợi ích du lịch từ đỉnh \(u\) tới đỉnh \(v\) là
\[
  \sum_i V_i W_{C_i},
\]
trong đó \(i\) chạy qua các màu, \(C_i\) là số lần màu \(i\) xuất hiện trên
đường đi đơn từ \(u\) tới \(v\), và \(W_0,W_1,\ldots,W_N\) là dãy số nguyên.

Cần thực hiện \(Q\) thao tác thuộc hai loại:
\begin{itemize}
  \item \texttt{CHANGE x\_i y\_i}: đổi màu của đỉnh \(x_i\) thành \(y_i\).
  \item \texttt{QUERY u\_i v\_i}: hỏi lợi ích du lịch từ \(u_i\) tới
  \(v_i\).
\end{itemize}

Quy mô dữ liệu: \(1\le N,M,Q\le 10^5\), \(1\le V_i,W_i\le 10^6\). Giới hạn
thời gian của bài là \(10\) giây. Bài gốc là \emph{Candy Park} của NOI
Winter Camp 2013.

\paragraph{Phân tích thuật toán.}
Chọn một đỉnh bất kỳ làm gốc. Đặt \(S=\Theta(N^{1/3})\), dùng phương pháp
đánh dấu đỉnh then chốt để tiền xử lý trên cây, ghi số lần xuất hiện của các
màu \(1\) tới \(M\) và lợi ích du lịch trên đường đi giữa mọi cặp đỉnh then
chốt. Độ phức tạp tiền xử lý là \(O(N^{5/3})\).

Với thao tác truy vấn đường đi từ \(u_i\) tới \(v_i\), trước hết dùng nhân đôi
hoặc thuật toán hiệu quả khác để tìm LCA của hai đỉnh. Sau đó lần lượt đi từ
\(u_i\) và \(v_i\) dọc theo đường đi để tìm các đỉnh then chốt gần \(u_i\) và
\(v_i\) nhất trên đường. Khi đó có thể dùng thông tin đã ghi và một số thao
tác tăng lượng để thu được đáp án trong \(O(N^{1/3})\) thời gian.

Với thao tác đổi màu đỉnh \(x_i\), liệt kê mọi cặp đỉnh then chốt và kiểm tra
\(x_i\) có nằm trên đường đi giữa chúng hay không. Nếu có thì cập nhật thông
tin đã ghi.

\paragraph{Bổ đề 3.13.}
Đỉnh \(x_i\) nằm trên đường đi từ \(u\) tới \(v\) khi và chỉ khi LCA của
\(u\) và \(v\) là tổ tiên của \(x_i\), đồng thời \(x_i\) là tổ tiên của
\(u\) hoặc \(v\).

Nếu tiền xử lý thời điểm vào stack và ra khỏi stack trong DFS cho mỗi đỉnh,
thì \(x\) là tổ tiên của \(y\) khi và chỉ khi thời điểm vào stack của \(y\)
nằm giữa thời điểm vào và ra của \(x\). Trong lúc tiền xử lý, ta ghi LCA của
mọi cặp đỉnh then chốt, nên có thể \(O(1)\) kiểm tra một đỉnh có nằm trên
đường đi giữa chúng hay không. Vì vậy thao tác sửa có độ phức tạp
\(O(N^{2/3})\).

\paragraph{Định lý 3.14.}
Bài toán này có thể giải trong \(O((N+Q)N^{2/3})\) thời gian.

\subsection{Chia khối và thuật toán ngoại tuyến}

Một số bài xử lý dữ liệu không yêu cầu trả lời từng truy vấn ngay lập tức, mà
cho trước rất nhiều truy vấn và yêu cầu chương trình nhanh chóng tính toàn bộ
kết quả rồi trả về một lần. Thuật toán ngoại tuyến là thuật toán biết trước
nội dung của mọi truy vấn, vì vậy có thể giải chúng cùng nhau.

Trong bài ngoại tuyến, vì đã biết toàn bộ truy vấn, ta có thể tổ chức các
truy vấn theo một cách nhất định, tận dụng đầy đủ thông tin trùng lặp giữa
chúng. Về bản chất, ta xem toàn bộ truy vấn như một phần của tập dữ liệu, rồi
dùng phương pháp xử lý tập dữ liệu để xử lý truy vấn.

Mục này dùng hai ví dụ để minh họa ứng dụng của tư tưởng chia khối khi thiết
kế thuật toán ngoại tuyến.

\paragraph{Ví dụ năm.}
Cho một dãy gồm \(N\) số, mỗi số là một số nguyên trong khoảng \(1\) tới
\(N\). Cần trả lời \(M\) truy vấn. Mỗi truy vấn cho ba số nguyên
\(L_i,R_i,K_i\), hỏi trong đoạn \([L_i,R_i]\) có bao nhiêu số xuất hiện đúng
\(K_i\) lần.

Quy mô dữ liệu: \(1\le N,M\le 10^5\).

\paragraph{Phân tích bài toán.}
Chia dãy thành khối, đặt kích thước mỗi khối là \(O(\sqrt N)\). Số khối cũng
là \(O(\sqrt N)\).

Đọc toàn bộ truy vấn một lần và phân loại truy vấn theo khối chứa đầu trái
\(L_i\), sau đó xử lý riêng từng loại.

Xét mọi truy vấn có đầu trái nằm trong khối \(A_i\). Sắp xếp các truy vấn này
theo vị trí đầu phải trong dãy từ nhỏ tới lớn. Sau đó quét lần lượt mọi phần
tử từ đầu trái của khối \(A_i\) tới cuối dãy, dùng mảng duy trì tại từng thời
điểm số lượng các giá trị có tần suất \(1\) tới \(N\) trong các phần tử đã
quét, và tần suất của từng giá trị từ \(1\) tới \(N\).

Với mỗi truy vấn \([L_i,R_i]\), khi quét tới đầu phải \(R_i\), ta có thể
trong \(O(\sqrt N)\) thời gian chuyển thông tin đang duy trì thành thông tin
của đoạn \([L_i,R_i]\). Sau khi ghi đáp án của truy vấn, khôi phục thông tin.
Mỗi truy vấn được xử lý đúng một lần.

\paragraph{Định lý 4.1.}
Thuật toán này có độ phức tạp thời gian \(O((N+M)\sqrt N)\) và độ phức tạp
không gian \(O(N)\).

Trong ví dụ này, ta dùng cách ngoại tuyến để giải quyết vấn đề dung lượng ghi
thông tin quá lớn của ví dụ ba. Có thể thấy kết hợp chia khối với thuật toán
ngoại tuyến có phần tương tự phương pháp tiền xử lý; hơn nữa, những bài có
thể giải bằng chia khối và tiền xử lý thường có thể giảm độ phức tạp không
gian trong tình huống ngoại tuyến.

\paragraph{Ví dụ sáu.}
Cho một cây \(N\) đỉnh. Mỗi đỉnh có trọng số là một số nguyên không âm nhỏ hơn
\(2^{31}\). Cần trả lời \(M\) truy vấn; mỗi truy vấn hỏi số lượng trọng số
đỉnh phân biệt xuất hiện trên đường đi đơn từ đỉnh \(u_i\) tới đỉnh \(v_i\).

Quy mô dữ liệu: \(1\le N\le 40000\), \(1\le M\le 100000\). Bài gốc là
\emph{Count On A Tree II} trên Sphere Online Judge.

\paragraph{Phân tích bài toán.}
Trước hết rời rạc hóa mọi trọng số của đỉnh; sau đó có thể xem trọng số là
số nguyên không vượt quá \(N\). Dùng phương pháp đánh dấu đỉnh then chốt trên
cây đã giới thiệu ở mục 3.3. Đặt \(S=\Theta(\sqrt N)\); khi đó số đỉnh then
chốt cũng là \(O(\sqrt N)\). Theo hệ quả 3.11, khoảng cách từ mỗi đỉnh tới
đỉnh then chốt gần nhất là \(O(\sqrt N)\).

Đọc toàn bộ truy vấn và phân loại mỗi truy vấn theo đỉnh then chốt gần nhất
với đầu mút \(u_i\) tương ứng. Sau đó xử lý riêng từng loại.

Xét một loại truy vấn có chung đỉnh then chốt \(r_i\), là đỉnh then chốt gần
nhất với các đầu mút \(u_i\). Lấy \(r_i\) làm gốc và DFS toàn bộ cây; đồng
thời duy trì, trong các đỉnh đã vào stack nhưng chưa ra stack, số lần xuất
hiện của từng trọng số và số lượng trọng số phân biệt.

\paragraph{Bổ đề 4.2.}
Khi đang thăm một đỉnh \(x_i\), mọi đỉnh đã vào stack nhưng chưa ra stack
chính là toàn bộ đỉnh trên đường từ \(x_i\) tới gốc.

\paragraph{Bổ đề 4.3.}
Đường đi từ một đỉnh \(u\) tới một đỉnh \(v\) tương đương với đường từ \(u\)
tới gốc cộng với đường từ \(v\) tới gốc, rồi trừ đường từ
\(\operatorname{LCA}(u,v)\) tới gốc.

Với mỗi truy vấn \((u_i,v_i)\), khi DFS tới \(v_i\), thông tin trong stack
chính là thông tin của các đỉnh trên đường từ \(v_i\) tới gốc. Lúc này liệt
kê từng đỉnh \(x\) trên đường từ \(u_i\) tới gốc \(r_i\). Nếu \(x\) là
\(\operatorname{LCA}(u_i,v_i)\) thì không thay đổi; nếu \(x\) nằm trên đường
từ \(\operatorname{LCA}(u_i,v_i)\) tới gốc \(r_i\) thì trừ trọng số của
\(x\); ngược lại cộng trọng số của \(x\). Sau khi thu được đáp án, hoàn tác
các thay đổi tạm thời.

\paragraph{Bổ đề 4.4.}
Có thể trả lời một truy vấn trong \(O(\sqrt N)\) thời gian.

\paragraph{Định lý 4.5.}
Bài toán này có thể giải trong \(O((N+M)\sqrt N)\) thời gian và \(O(N)\)
không gian.

\subsection{Tổng kết}

Nói một cách hình tượng, bản chất của tư tưởng chia khối là chuyển một thuật
toán đơn giản có độ phức tạp \(O(f(N))\) thành một thuật toán có độ phức tạp
\[
  O\left(g(S)+\frac{N}{S}h(S)\right),
\]
trong đó \(S\) là kích thước khối, còn \(N/S\) là số khối. Với các bài khác
nhau, biểu hiện cụ thể của điều này cũng khác nhau.

Lấy ví dụ một làm mẫu: thuật toán tham khảo của bài đó thể hiện tư tưởng chia
khối ở hai chỗ. Một chỗ là dùng hai cấu trúc dữ liệu khác nhau để duy trì các
sửa có \(D_i<\sqrt N\) và các sửa có \(D_i\ge\sqrt N\). Chỗ còn lại là khi
thao tác sửa rất thường xuyên, ta dùng cấu trúc chia khối để duy trì tổng
đoạn của dãy. Trong bài đặc biệt này, cách đó hiệu quả hơn cây đoạn, vì số
lượng thao tác sửa có thể đạt tới \(O(\sqrt N)\) lần số lượng truy vấn. Các
ví dụ sau cũng tương tự. Tuy nhiên, có khi ta chia một bài toán thành nhiều
bài con có bản chất giống nhau, có khi lại tách một bài toán thành nhiều bài
con khác nhau khá lớn.

Bài viết đưa ra các mô hình cơ bản của chia khối trên dãy tuyến tính và cấu
trúc cây; chúng có hiệu lực với một lớp bài rất rộng. Nhưng áp dụng tư tưởng
chia khối như thế nào còn tùy từng bài cụ thể. Các bạn không nên bị bó buộc
bởi vài hình thức được nêu ở đây, mà cần mạnh dạn sáng tạo. Tất nhiên, nếu
tồn tại thuật toán hiệu quả hơn chia khối, cũng phải biết lựa chọn, không nên
bảo thủ. Hy vọng các ví dụ này đem lại gợi ý cho người đọc.

\subsection{Tài liệu tham khảo}

\begin{enumerate}
  \item Thomas H. Cormen, Charles E. Leiserson, Ronald L. Rivest, Clifford
  Stein. \emph{Introduction to Algorithms}, 2nd edition.
  \item Jon Kleinberg, Eva Tardos. \emph{Algorithm Design}.
  \item Wu Wenhu, Wang Jiande. \emph{Giáo trình nâng cao cuộc thi lập trình
  sinh viên quốc tế ACM/ICPC}, tập 1.
  \item Zheng Dun. \emph{Quy hoạch cân bằng: phân tích ứng dụng của một lớp
  tư tưởng cân bằng}.
  \item Qi Zichao. \emph{Ứng dụng thuật toán chia để trị trong các bài toán
  đường đi trên cây}.
  \item Tài liệu về \emph{Candy Park}, NOI Winter Camp 2013.
  \item Báo cáo ra đề \emph{JZPLCM}.
  \item Xu Jie. Báo cáo ra đề \emph{Đồ thị trong đồ thị}.
  \item Xu Haoran. \emph{Bàn rộng về cấu trúc dữ liệu}.
\end{enumerate}

\subsection{Lời cảm ơn}

Xin cảm ơn China Computer Federation đã cung cấp nền tảng học tập và trao đổi.

Xin cảm ơn thầy hướng dẫn Gong Zhiyong đã quan tâm và giúp đỡ tác giả trong
nhiều năm.

Xin cảm ơn các huấn luyện viên đội tuyển quốc gia Hu Weidong và Tang Wenbin
vì sự giúp đỡ.

Đặc biệt cảm ơn bạn Fan Haoqiang ở Bắc Kinh đã kiên nhẫn giúp đỡ tác giả
trong quá trình viết bài. Tác giả học được rất nhiều và nhận được không ít
gợi ý qua các trao đổi với bạn ấy.

Xin cảm ơn đàn anh Huang Jiyuan vì sự giúp đỡ.

Xin cảm ơn các thầy cô và bạn học khác từng giúp đỡ hoặc gợi ý cho tác giả.

Cuối cùng, xin cảm ơn cha mẹ tôi vì sự quan tâm và chăm sóc chu đáo.

\section{Kỹ thuật \emph{meet in the middle} trong bài toán tìm kiếm}

\begin{center}
Qiao Mingda\\
Trường Ngoại ngữ Nam Kinh, Giang Tô
\end{center}

\subsection{Tóm tắt}

\emph{Meet in the middle} là một kỹ thuật tối ưu hóa tìm kiếm thường gặp. Một
số bài toán tìm kiếm có thể dùng kỹ thuật này để giảm mạnh độ phức tạp thời
gian, nhờ đó tìm được kết quả trong giới hạn thời gian.

Bài viết gồm ba phần. Phần thứ nhất đưa ra một mô hình đồ thị có hướng trừu
tượng, dựa vào đó giới thiệu tư tưởng của kỹ thuật \emph{meet in the middle},
và dùng kỹ thuật này để thiết kế thuật toán cho mô hình đồ thị có hướng. Phần
thứ hai chứng minh ngắn gọn tính đúng đắn, độ phức tạp thời gian và điều kiện
sử dụng của thuật toán trong phần thứ nhất. Phần thứ ba khảo sát vài bài ví
dụ dùng kỹ thuật này; từ các ví dụ đó khái quát thành một mô hình phương
trình, rồi phân tích một số biến thể của kỹ thuật.

\subsection{Phần thứ nhất: mô hình đồ thị có hướng}

\subsubsection{Bản chất bài toán}

Bản chất của một phần các bài toán tìm kiếm là: trong đồ thị có hướng \(G\),
tính số đường đi khác nhau có độ dài \(L\) từ đỉnh \(A\) tới đỉnh \(B\). Ở
đây, một ``đường đi độ dài \(L\)'' được định nghĩa là một dãy cạnh có hướng
\[
  (E_1,E_2,\ldots,E_L),
\]
trong đó điểm đầu của \(E_1\) là \(A\), điểm cuối của \(E_L\) là \(B\), và
với mọi \(i=1,2,\ldots,L-1\), điểm cuối của \(E_i\) bằng điểm đầu của
\(E_{i+1}\). Hai đường đi
\[
  (E_1,E_2,\ldots,E_L)
  \quad\text{và}\quad
  (E'_1,E'_2,\ldots,E'_L)
\]
khác nhau khi và chỉ khi tồn tại \(1\le i\le L\) sao cho \(E_i\ne E'_i\).

Giả sử bậc ra lớn nhất của một đỉnh trong \(G\) là hằng số \(D\), không gian
cần để mô tả một đỉnh là \(M\), và thời gian cần để tính một hậu kế xác định
của một đỉnh là \(T\).

Ví dụ, mỗi đỉnh có nhãn là một dãy độ dài \(N\),
\[
  P_1,P_2,\ldots,P_N.
\]
Từ một đỉnh, ta có thể lật dãy \(P_1,P_2,\ldots,P_{N-1}\) hoặc dãy
\(P_2,P_3,\ldots,P_N\) để thu được hậu kế. Ở đây, ``lật
\(P_1,P_2,\ldots,P_{N-1}\)'' nghĩa là đổi chỗ \(P_1\) với \(P_{N-1}\),
\(P_2\) với \(P_{N-2}\), v.v., cho tới \(P_{\lfloor N/2\rfloor}\) và
\(P_{\lceil N/2\rceil}\). Như vậy mỗi đỉnh có đúng hai hậu kế; mô tả một
đỉnh cần \(O(N)\) không gian để lưu dãy; tính một hậu kế cần \(O(N)\) thời
gian để lật dãy. Trong ví dụ này \(D=2\), \(M=O(N)\), và \(T=O(N)\).

Theo các giả thiết trên, ta phân tích thuật toán DFS thô để giải bài toán
này. Thuật toán giới hạn độ sâu DFS không vượt quá \(L\); khi độ sâu tìm kiếm
đạt \(L\), nó so sánh đỉnh hiện tại với \(B\), nếu là cùng một đỉnh thì tăng
đáp án thêm \(1\). Dễ chứng minh số đỉnh đi qua trong quá trình tìm kiếm
không vượt quá
\[
  1+D+D^2+\cdots+D^L=\frac{D^{L+1}-1}{D-1}=O(D^L),
\]
và số đỉnh đạt được sau khi tìm \(L\) bước không vượt quá \(D^L=O(D^L)\). Lưu
ý rằng nếu cùng một đỉnh được đi qua nhiều lần thì phải tính nhiều lần.

Vì một đỉnh chiếm \(O(M)\) không gian, so sánh một đỉnh có phải \(B\) hay
không cần \(O(M)\) thời gian. Mặt khác, trong quá trình tìm kiếm, tính mỗi
hậu kế cần \(O(T)\) thời gian. Vì vậy độ phức tạp thời gian của DFS thô là
\[
  O(D^L(M+T)).
\]
Khi \(L\) lớn, thời gian chạy của thuật toán thô sẽ rất dài; trong đa số tình
huống nó không thể tính ra đáp án trong giới hạn thời gian. Ta cần xét cách
tối ưu hóa tìm kiếm.

Trong lớp bài này, số đỉnh của \(G\) rất lớn, thậm chí \(G\) có thể là đồ thị
vô hạn, nên số bài toán con trùng lặp trong quá trình tìm kiếm khá ít; quy
hoạch động không phải một phương pháp hiệu quả. Vì trong trường hợp xấu nhất
ta buộc phải duyệt cả cây tìm kiếm, các kỹ thuật tối ưu hóa tìm kiếm nghiệm
tối ưu cũng không phù hợp với loại bài đếm số đường đi này. Ta cần thay đổi
cách tìm kiếm ở mức tổng thể.

\subsubsection{Tư tưởng \emph{meet in the middle}}

Nghĩa đen của \emph{meet in the middle} là ``gặp nhau ở giữa''. Giả sử có hai
người, một người xuất phát từ \(A\), người kia xuất phát từ \(B\), cùng đi về
phía điểm xuất phát của người còn lại. Nếu khi cả hai đều đi được \(L/2\) mét
họ vừa đúng ``gặp nhau ở giữa'', thì có thể khẳng định ta đã tìm được một
đường đi từ \(A\) tới \(B\) có độ dài \(L\) mét.

Dưới đây ta xét cách áp dụng tư tưởng này vào mô hình đồ thị có hướng.

\subsubsection{Quy trình thuật toán cho mô hình đồ thị có hướng}

Để thuận tiện mô tả, ta dựng một đồ thị mới \(G'\). Các đỉnh của \(G'\) giống
với \(G\), nhưng mỗi cạnh có hướng từ \(X\) tới \(Y\) trong \(G\) tương ứng
với một cạnh có hướng từ \(Y\) tới \(X\) trong \(G'\). Ngoài ra, đặt
\[
  L_1=\left\lfloor\frac{L}{2}\right\rfloor,\qquad
  L_2=\left\lceil\frac{L}{2}\right\rceil.
\]
Xét riêng tính chẵn lẻ của \(L\) sẽ dễ chứng minh \(L_1+L_2=L\).

Phần thứ nhất của thuật toán thực hiện DFS trong \(G\) từ đỉnh \(A\), giới
hạn độ sâu tìm kiếm không vượt quá \(L_1\). Ta ghi lại toàn bộ các đỉnh đạt
được sau đúng \(L_1\) bước và số lần đạt được từng đỉnh. Ký hiệu số lần đi
\(L_1\) bước tới đỉnh \(P\) là \(\operatorname{Count}(P)\). Nếu sau
\(L_1\) bước chưa từng tới \(P\), quy ước \(\operatorname{Count}(P)=0\).

Phần thứ hai của thuật toán thực hiện DFS trong \(G'\) từ đỉnh \(B\), giới
hạn độ sâu tìm kiếm không vượt quá \(L_2\). Nếu sau đúng \(L_2\) bước tìm
kiếm tới đỉnh \(Q\), ta tuyên bố đã tìm được
\(\operatorname{Count}(Q)\) đường đi có độ dài \(L_1+L_2=L\), và cộng
\(\operatorname{Count}(Q)\) vào đáp án. Lưu ý rằng sau \(L_2\) bước tìm kiếm
ta có thể tới cùng một đỉnh nhiều lần; mỗi lần tới đều phải cộng
\(\operatorname{Count}(Q)\).

Từ quy trình trên có thể thấy, \emph{meet in the middle} là một kỹ thuật
``tìm kiếm hai chiều''.

\subsection{Phần thứ hai: phân tích thuật toán}

\subsubsection{Tính đúng đắn}

Đặt
\[
  L_1=\left\lfloor\frac{L}{2}\right\rfloor,\qquad
  L_2=\left\lceil\frac{L}{2}\right\rceil.
\]
Với một đỉnh \(P\) trong \(G\), gọi \(f(P)\) là số đường đi từ \(A\) tới
\(P\) có độ dài \(L_1\), và gọi \(g(P)\) là số đường đi từ \(P\) tới \(B\)
có độ dài \(L_2\).

\paragraph{Bổ đề 1.}
Số đường đi từ \(A\) tới \(B\) có độ dài \(L\) là
\[
  \sum_P f(P)g(P).
\]

\paragraph{Chứng minh.}
Xét một đường đi bất kỳ có độ dài \(L\),
\[
  (E_1,\ldots,E_L),
\]
trong đó \(E_1\) bắt đầu ở \(A\), còn \(E_L\) kết thúc ở \(B\). Lấy \(P\) là
điểm cuối của \(E_{L_1}\); ta nói đường đi này ``đi qua'' \(P\).

Ngược lại, nếu chọn một đường đi từ \(A\) tới \(P\) độ dài \(L_1\), rồi chọn
một đường đi từ \(P\) tới \(B\) độ dài \(L_2\), ghép hai dãy cạnh theo thứ tự
thì thu được một đường đi từ \(A\) tới \(B\) độ dài \(L\). Vì vậy, với một
đỉnh \(P\) cố định, số đường đi ``đi qua'' \(P\) là \(f(P)g(P)\). Lấy tổng
trên mọi \(P\), ta được số đường đi cần tính.

\paragraph{Bổ đề 2.}
Thuật toán trong phần thứ nhất cho ra kết quả
\[
  \sum_P f(P)g(P).
\]

\paragraph{Chứng minh.}
Theo hai lần DFS, với mọi đỉnh \(P\), trong phần thứ nhất của thuật toán ta
có \(\operatorname{Count}(P)=f(P)\); trong phần thứ hai, số lần tìm kiếm sau
\(L_2\) bước đạt tới \(P\) là \(g(P)\).

Do đó, với các đỉnh có \(g(P)=0\), thuật toán không tăng đáp án, tương đương
cộng thêm \(f(P)g(P)=0\). Với các đỉnh có \(g(P)>0\), thuật toán cộng
\(f(P)\) vào đáp án đúng \(g(P)\) lần, tức cộng tổng cộng \(f(P)g(P)\). Cộng
theo mọi đỉnh \(P\), kết quả của thuật toán là \(\sum_P f(P)g(P)\).

\paragraph{Định lý 1.}
Thuật toán trong phần thứ nhất có thể tính đúng tổng số đường đi khác nhau từ
đỉnh \(A\) tới đỉnh \(B\) có độ dài \(L\) trong đồ thị \(G\).

\paragraph{Chứng minh.}
Kết hợp Bổ đề 1 và Bổ đề 2, đáp án thuật toán đưa ra và số đường đi thực tế
đều bằng \(\sum_P f(P)g(P)\).

\subsubsection{Độ phức tạp thời gian}

Dưới đây ta phân tích độ phức tạp thời gian của thuật toán. Ở trên ta đã giả
sử bậc ra lớn nhất của một đỉnh trong \(G\) là \(D\); bây giờ giả sử thêm
rằng bậc ra của các đỉnh trong \(G'\) cũng không vượt quá \(D\). Mặt khác,
giống mô tả trước đó, ta vẫn giả sử trong cả \(G\) và \(G'\), không gian mô
tả một đỉnh là \(M\), và thời gian tính hậu kế xác định của một đỉnh là
\(T\).

Độ phức tạp thời gian của thuật toán chủ yếu phụ thuộc vào hai phần: thời
gian DFS và thời gian chương trình tính, truy vấn \(\operatorname{Count}(P)\).
Trong phân tích dưới đây, vì \(D\) là hằng số, ta có thể bỏ qua các dấu làm
tròn trên số mũ, tức xem
\[
  O(D^{\lfloor L/2\rfloor})
  =O(D^{\lceil L/2\rceil})
  =O(D^{L/2}).
\]

Trong phần thứ nhất, chi phí DFS là \(O(D^{L/2}T)\); trong phần thứ hai, chi
phí DFS cũng là \(O(D^{L/2}T)\). Vì vậy tổng thời gian dùng cho DFS là
\[
  O(D^{L/2}T)+O(D^{L/2}T)=O(D^{L/2}T).
\]

Trong phần thứ nhất, ta cần thống kê số lần
\(\operatorname{Count}(P)\) của mỗi đỉnh \(P\) đạt được sau \(L_1\) bước; trong
phần thứ hai, ta cần truy vấn \(\operatorname{Count}(Q)\) với mỗi đỉnh \(Q\)
đạt được sau \(L_2\) bước. Điều này đòi hỏi một cấu trúc dữ liệu có thể chèn
dữ liệu nhanh và truy vấn nhanh số lần xuất hiện của một dữ liệu. Có thể thấy
bảng băm là lựa chọn lý tưởng. Theo tính chất của bảng băm, nếu dùng băm đều
đơn giản (\emph{simple uniform hashing}) và dùng danh sách liên kết để xử lý
va chạm, thì trung bình một thao tác chèn hoặc truy vấn chỉ cần \(O(M)\) thời
gian. Nguyên nhân là trong bài toán này, so sánh hai đỉnh có giống nhau hay
không cần \(O(M)\) thời gian.

Trong phần thứ nhất, ta cần thực hiện \(O(D^{L/2})\) thao tác chèn vào bảng
băm; trong phần thứ hai, ta cần thực hiện \(O(D^{L/2})\) thao tác truy vấn.
Vì vậy thời gian tính và truy vấn \(\operatorname{Count}(P)\) là
\[
  \bigl[O(D^{L/2})+O(D^{L/2})\bigr]\cdot O(M)
  =O(D^{L/2}M).
\]

Kết hợp hai phần, độ phức tạp thời gian của thuật toán là
\[
  O(D^{L/2}(M+T)).
\]

\subsubsection{Điều kiện sử dụng thuật toán}

Muốn dùng thuật toán trong phần thứ nhất cho một bài toán, trước hết bài toán
đó phải phù hợp với mô hình đồ thị có hướng. Tiếp theo, trong quá trình phân
tích độ phức tạp thời gian của thuật toán, thực ra ta đã thêm một số điều
kiện hạn chế.

Trong phân tích trước đó, ta giả sử bậc ra của các đỉnh trong \(G'\) không
vượt quá bậc ra lớn nhất trong \(G\). Mặt khác, bậc ra của một đỉnh trong
\(G'\) thực chất bằng bậc vào của đỉnh đó trong \(G\). Vì vậy thông thường
cần bảo đảm bậc vào và bậc ra của các đỉnh trong \(G\) cùng bậc độ lớn.

Mặt khác, ta giả sử có thể truy cập một hậu kế xác định của một đỉnh trong
\(G'\) trong \(O(T)\) thời gian. Điều kiện này tương đương với yêu cầu có thể
truy cập một tiền nhiệm xác định của một đỉnh trong \(G\) trong \(O(T)\) thời
gian.

Về chi phí không gian, để bảo đảm bảng băm có độ phức tạp thao tác kỳ vọng
\(O(M)\), ta cần tiêu tốn
\[
  O(D^{L/2}M)
\]
không gian.

\subsection{Phần thứ ba: ứng dụng}

\subsubsection{Ví dụ 1: số nghiệm của phương trình}

\paragraph{Nội dung bài toán.}
Cho các số nguyên \(M\), \(k_1,k_2,\ldots,k_N\) và
\(p_1,p_2,\ldots,p_N\). Hãy tính số nghiệm nguyên của phương trình
\[
  \sum_{i=1}^{N} k_i x_i^{p_i}=0
\]
thỏa mãn \(1\le x_i\le M\). Bài gốc là NOI 2001, vòng hai.

\paragraph{Quy mô dữ liệu.}
\[
  N\le 6,\qquad M\le 150.
\]

\paragraph{Phân tích.}
Ta có thể chuyển bài này thành mô hình đồ thị có hướng, rồi dùng thuật toán
trong phần thứ nhất để giải.

Dùng cặp \((\operatorname{Depth},\operatorname{Sum})\) để biểu diễn một đỉnh
trong \(G\), nên mô tả một đỉnh chỉ cần \(O(1)\) không gian. Ở đây
\(\operatorname{Depth}\) biểu diễn số hạng đã xử lý, còn
\(\operatorname{Sum}\) biểu diễn tổng của \(\operatorname{Depth}\) số hạng
đầu. Đỉnh xuất phát là \((0,0)\), đỉnh kết thúc là \((N,0)\).

Với một đỉnh \((\operatorname{Depth},\operatorname{Sum})\) có
\(\operatorname{Depth}<N\), đặt
\[
  k'=k_{\operatorname{Depth}+1},\qquad
  p'=p_{\operatorname{Depth}+1}.
\]
Từ đỉnh này có \(M\) hậu kế:
\[
\begin{aligned}
  &(\operatorname{Depth}+1,\operatorname{Sum}+k'\cdot 1^{p'}),\\
  &(\operatorname{Depth}+1,\operatorname{Sum}+k'\cdot 2^{p'}),\\
  &\ldots,\\
  &(\operatorname{Depth}+1,\operatorname{Sum}+k'\cdot M^{p'}).
\end{aligned}
\]

Nếu xét đồ thị đảo, thì với một đỉnh \((\operatorname{Depth},\operatorname{Sum})\)
có \(\operatorname{Depth}>0\), nó có \(M\) tiền nhiệm. Đặt
\[
  k''=k_{\operatorname{Depth}},\qquad p''=p_{\operatorname{Depth}}.
\]
Các tiền nhiệm là
\[
\begin{aligned}
  &(\operatorname{Depth}-1,\operatorname{Sum}-k''\cdot 1^{p''}),\\
  &(\operatorname{Depth}-1,\operatorname{Sum}-k''\cdot 2^{p''}),\\
  &\ldots,\\
  &(\operatorname{Depth}-1,\operatorname{Sum}-k''\cdot M^{p''}).
\end{aligned}
\]

Để tính một hậu kế hoặc tiền nhiệm của một đỉnh, ta cần tính một lũy thừa.
Dùng lũy thừa nhanh thì độ phức tạp thời gian là \(O(\log \operatorname{MaxP})\),
trong đó \(\operatorname{MaxP}\) là giá trị lớn nhất trong
\(p_1,p_2,\ldots,p_N\).

Qua phân tích trên, ta có thể dùng trực tiếp thuật toán trong phần thứ nhất,
với độ phức tạp thời gian
\[
  O(M^{N/2}\log \operatorname{MaxP}).
\]

Hãy nhìn lại ý nghĩa thực tế của thuật toán trên trong ví dụ này. Xét trường
hợp \(N=6\):
\[
  k_1x_1^{p_1}+k_2x_2^{p_2}+k_3x_3^{p_3}
  +k_4x_4^{p_4}+k_5x_5^{p_5}+k_6x_6^{p_6}=0.
\]
Thuật toán thô là liệt kê sáu ẩn ở vế trái rồi kiểm tra kết quả của vế trái
có bằng hằng số ở vế phải hay không; độ phức tạp thời gian như vậy là
\(O(M^6\log \operatorname{MaxP})\), rất kém.

Suy nghĩ kỹ hơn, phương trình trên có vẻ rất mất cân bằng: cả sáu ẩn đều nằm
ở vế trái, khiến ta buộc phải liệt kê sáu ẩn. Chuyển vế phương trình, ta có
\[
  k_1x_1^{p_1}+k_2x_2^{p_2}+k_3x_3^{p_3}
  =
  -k_4x_4^{p_4}-k_5x_5^{p_5}-k_6x_6^{p_6}.
\]
Lúc này hai vế đã ``cân bằng''. Giả sử ta liệt kê trước ba số ở vế trái và
lưu mọi kết quả liệt kê được vào một bảng, rồi liệt kê ba số ở vế phải. Mỗi
khi tính được một kết quả trùng với một số trong bảng, ta thu được một
nghiệm của phương trình. Khi so sánh kết quả tính toán, ta không cần so sánh
từng cặp số; có thể dùng cấu trúc dữ liệu, chẳng hạn bảng băm, để tăng tốc
việc thống kê. Từ đó ta có thuật toán độ phức tạp
\(O(M^3\log \operatorname{MaxP})\). Thực chất, điểm then chốt của
\emph{meet in the middle} nằm ở việc giảm lượng tính toán bằng cách tối ưu
hóa quá trình hợp nhất.

\subsubsection{Ví dụ 2: \emph{ABCDEF}}

\paragraph{Nội dung bài toán.}
Cho một tập số nguyên \(S\). Hãy tính có bao nhiêu bộ
\((a,b,c,d,e,f)\) thỏa mãn \(a,b,c,d,e,f\in S\), \(d\ne 0\), và
\[
  ab+c=d(e+f).
\]
Bài gốc là \emph{ABCDEF} trên SPOJ.

\paragraph{Quy mô dữ liệu.}
\[
  |S|\le 100.
\]

\paragraph{Phân tích.}
Bài này chỉ cần chuyển vế phương trình thành
\[
  ab+c=d(e+f),
\]
vì vậy về bản chất nó giống ví dụ 1. Ta thu được thuật toán có độ phức tạp
\[
  O(|S|^3).
\]

\subsubsection{Ví dụ 3: \emph{EllysBulls}}

\paragraph{Nội dung bài toán.}
Có một số thập phân \(N\) chữ số \(A\), có thể có chữ số \(0\) ở đầu. Bây giờ
cho \(M\) điều kiện; điều kiện thứ \(i\) gồm một số \(N\) chữ số \(X_i\) và
giá trị \(\operatorname{Same}(A,X_i)\). Ở đây \(\operatorname{Same}(A,B)\)
biểu diễn số vị trí mà hai số \(A\) và \(B\) có chữ số giống nhau. Ví dụ,
\[
  \operatorname{Same}(\text{``0312''},\text{``0321''})=2.
\]
Nếu có duy nhất một \(A\) thỏa mãn yêu cầu thì in nghiệm đó; ngược lại in
thông tin biểu thị vô nghiệm hoặc nhiều nghiệm. Bài gốc là TopCoder Single
Round Match 572, Division I, Level Two.

\paragraph{Quy mô dữ liệu.}
\[
  N\le 9,\qquad M\le 50.
\]

\paragraph{Phân tích.}
Để tiện trình bày, định nghĩa vector hàng
\[
  \operatorname{Sum}
  =
  \bigl(\operatorname{Same}(A,X_1),\operatorname{Same}(A,X_2),
  \ldots,\operatorname{Same}(A,X_M)\bigr).
\]
Gọi \(X_{i,j}\) là chữ số thứ \(j\) của \(X_i\). Định nghĩa hàm \(S(a,b)\)
với \(0\le a,b\le 9\): nếu \(a=b\) thì \(S(a,b)=1\), nếu \(a\ne b\) thì
\(S(a,b)=0\). Tiếp tục định nghĩa
\[
  \operatorname{Cost}(i,j)
  =
  \bigl(S(X_{1,i},j),S(X_{2,i},j),\ldots,S(X_{M,i},j)\bigr).
\]
Vector này cho biết: nếu chữ số thứ \(i\) của \(A\) là \(j\), thì riêng vị
trí này khớp với những điều kiện nào.

Theo điều kiện của đề bài, ta có phương trình
\[
  \operatorname{Cost}(1,a_1)+\operatorname{Cost}(2,a_2)+\cdots+
  \operatorname{Cost}(N,a_N)=\operatorname{Sum}.
\]
Ta vẫn chuyển vế giống hai bài trước. Đặt
\[
  \operatorname{Mid}=\left\lfloor\frac{N}{2}\right\rfloor,
\]
thu được
\[
\begin{aligned}
  &\operatorname{Cost}(1,a_1)+\cdots+\operatorname{Cost}(\operatorname{Mid},a_{\operatorname{Mid}})\\
  &\quad=
  \operatorname{Sum}
  -\operatorname{Cost}(\operatorname{Mid}+1,a_{\operatorname{Mid}+1})
  -\cdots-\operatorname{Cost}(N,a_N).
\end{aligned}
\]
Vì mỗi vị trí có \(10\) khả năng, nên tìm kiếm vế trái có
\(O(10^{N/2})\) kết quả khác nhau, và vế phải cũng có \(O(10^{N/2})\) kết
quả khác nhau. Hai vế của phương trình đều là vector hàng \(1\times M\), nên
mỗi bước tính toán và thao tác trong bảng băm cần \(O(M)\) thời gian. Vì vậy
ta có thuật toán độ phức tạp \(O(10^{N/2}M)\), đủ để vượt qua dữ liệu kiểm
thử của bài này.

\subsubsection{Tiểu kết: mô hình phương trình}

Tổng kết ba ví dụ trên, ta thấy các bài toán đều có thể chuyển thành một mô
hình phương trình:
\[
  \sum_{i=1}^{N} f(i,x_i)=S.
\]
Ở ví dụ 1 và 2, \(f(i,x_i)\) và \(S\) là số nguyên; ở ví dụ 3,
\(f(i,x_i)\) và \(S\) là vector. Trong cả ba bài, cách làm đều dựa trên việc
biến đổi phương trình thành
\[
  \sum_{i=1}^{\lfloor N/2\rfloor} f(i,x_i)
  =
  S-\sum_{i=\lfloor N/2\rfloor+1}^{N} f(i,x_i).
\]
Ta tìm kiếm vế trái trước và ghi kết quả vào bảng băm; sau đó tìm kiếm vế
phải và truy vấn các kết quả đã lưu trong bảng băm. Có thể thấy quy trình
giải ngắn gọn này rất giống thuật toán trong phần thứ nhất. Hơn nữa, từ phân
tích của ví dụ 1, mô hình phương trình có thể chuyển thành mô hình đồ thị có
hướng trong phần thứ nhất, nên mô hình phương trình là một trường hợp đặc
biệt của mô hình đồ thị có hướng. Với các bài thuộc mô hình phương trình, ta
dễ nghĩ tới kỹ thuật \emph{meet in the middle}.

Vận dụng hai mô hình trên, ta có thể nhanh chóng mô hình hóa bài toán cụ thể
và dùng kỹ thuật \emph{meet in the middle} để giải. Tuy nhiên, một số bài
không thể chuyển thành mô hình phương trình hoặc mô hình đồ thị có hướng, do
đó không thể dùng trực tiếp thuật toán của phần thứ nhất. Dù vậy, chúng vẫn
có điểm tương tự với hai mô hình trên; ta vẫn có thể sửa thuật toán trong
phần thứ nhất để dùng kỹ thuật \emph{meet in the middle} tối ưu các bài này.
Dưới đây ta phân tích vài ví dụ như vậy.

\subsubsection{Ví dụ 4: \emph{Balanced Cow Subsets}}

\paragraph{Nội dung bài toán.}
Cho các số nguyên dương \(a_1,\ldots,a_N\). Trước hết chọn ra một số số trong
\(N\) số, ít nhất một số; sau đó chia các số đã chọn thành hai phần sao cho
tổng của hai phần bằng nhau. Hãy tính số phương án chọn số ở bước đầu.

Lưu ý: có thể tồn tại một phương án chọn số cho phép nhiều cách chia khác
nhau, nhưng phương án chọn số đó chỉ được tính một lần. Bài gốc là USACO 2012
US Open, Gold Division.

\paragraph{Quy mô dữ liệu.}
\[
  N\le 20.
\]

\paragraph{Phân tích.}
Ta đổi bài toán sang một cách mô tả tương đương: với dãy
\(a_1,a_2,\ldots,a_N\), đặt \(x_i\in\{-1,0,1\}\). Dưới điều kiện
\[
  \sum_{i=1}^{N} x_i a_i=0,
\]
cần tính có bao nhiêu tập khác nhau, không rỗng,
\[
  S=\{i\mid 1\le i\le N,\ x_i\ne 0\}.
\]

Bài này giống mô hình phương trình đã nêu, nhưng điều ta cần không phải số
nghiệm của phương trình, mà là số tập tương ứng với các nghiệm của phương
trình. Vì vậy thuật toán cần sửa đôi chút.

Định nghĩa một dãy mới \(y_1,y_2,\ldots,y_N\): nếu \(x_i=0\) thì \(y_i=0\),
còn nếu \(x_i\ne 0\) thì \(y_i=2^{i-1}\). Khi đó tập \(S\) có thể biểu diễn
bằng \(\sum y_i\).

Chọn một số nguyên \(M\) trong đoạn \([0,N]\), biến đổi
\(\sum x_i a_i=0\) thành
\[
  x_1a_1+x_2a_2+\cdots+x_Ma_M
  =
  -x_{M+1}a_{M+1}-x_{M+2}a_{M+2}-\cdots-x_Na_N.
\]
Ký hiệu vế trái là \(\operatorname{LeftSum}\), vế phải là
\(\operatorname{RightSum}\). Tương tự, ký hiệu
\[
  \operatorname{LeftSet}=y_1+y_2+\cdots+y_M,\qquad
  \operatorname{RightSet}=y_{M+1}+y_{M+2}+\cdots+y_N.
\]

Trong phần thứ nhất của thuật toán, ta vẫn tìm kiếm \(M\) ẩn ở vế trái như
các bài trước, tổng cộng có \(3^M\) trường hợp. Ta không chỉ ghi
\(\operatorname{LeftSum}\), mà còn ghi cả \(\operatorname{LeftSet}\) tương
ứng.

Trong phần thứ hai của thuật toán, nếu phát hiện
\(\operatorname{RightSum}\) bằng một \(\operatorname{LeftSum}\) nào đó, thì
tập được biểu diễn bởi
\[
  \operatorname{LeftSet}+\operatorname{RightSet}
\]
là một tập hợp hợp lệ. Ta có thể dùng một mảng boolean đánh số từ \(0\) tới
\(2^N-1\) để ghi nhận các tập này.

Cuối cùng, quét toàn bộ mảng boolean một lần là thống kê được đáp án.

Dễ thấy phần thứ nhất có độ phức tạp thời gian \(O(3^M)\). Phần thứ hai liệt
kê \(O(3^{N-M})\) trường hợp, nhưng với mỗi trường hợp cần truy cập tất cả
\(\operatorname{LeftSet}\) tương ứng. Trong trường hợp xấu nhất, ta cần truy
cập \(O(2^M)\) \(\operatorname{LeftSet}\), nên độ phức tạp thời gian của
phần thứ hai là \(O(3^{N-M}2^M)\). Thời gian quét cuối cùng là \(O(2^N)\).

Tóm lại, độ phức tạp thời gian là
\[
  O(3^M+3^{N-M}2^M+2^N)=O(3^M+3^{N-M}2^M).
\]
Nếu chọn \(M=\lfloor N/2\rfloor\), ta thu được
\[
  O(6^{N/2})\approx O(2.45^N),
\]
đủ để vượt qua dữ liệu kiểm thử của bài này. Nhưng thực tế, nếu chọn
\[
  M=N\log_{9/2}3,
\]
ta có được độ phức tạp tốt hơn:
\[
  O\bigl((3^{\log_{9/2}3})^N\bigr)\approx O(2.23^N).
\]

Ví dụ 4 cho thấy khi dùng kỹ thuật \emph{meet in the middle}, ta có thể lưu
thêm một số thông tin khác trong bảng băm để thuận tiện cho việc giải bài.
Mặt khác, khi dùng kỹ thuật này, nếu chia số ẩn cần tìm kiếm thành hai phần
không đều nhau, đôi khi ta có thể thu được độ phức tạp thời gian tốt hơn.

\subsubsection{Ví dụ 5: \emph{AlphabetPaths}}

\paragraph{Nội dung bài toán.}
Cho một ma trận \(R\times C\). Mỗi ô của ma trận hoặc là ô trống, hoặc chứa
một số nguyên từ \(0\) tới \(20\). Hãy tính có bao nhiêu đường đi gồm \(21\)
ô thỏa mãn: hai ô kề nhau trên đường đi phải có một cạnh chung, và mỗi số từ
\(0\) tới \(20\) xuất hiện đúng một lần trên đường đi. Bài gốc là TopCoder
Single Round Match 523, Division I, Level Three.

\paragraph{Quy mô dữ liệu.}
\[
  R,C\le 21.
\]

\paragraph{Phân tích.}
Bài này nhìn qua không phải mô hình phương trình đã nêu, mà giống mô hình đồ
thị có hướng trong phần thứ nhất. Tuy nhiên, trong bài này điểm đầu và điểm
cuối của đường đi đều không xác định, mỗi loại có \(O(RC)\) cách chọn. Nếu
liệt kê cả điểm đầu và điểm cuối thì ta đã mất \(O(R^2C^2)\) lần liệt kê;
nhân thêm thời gian tìm kiếm thì hoàn toàn không thể vượt qua.

Ý tưởng vừa nêu là liệt kê điểm đầu và điểm cuối, rồi tìm từ hai đầu đường đi
vào giữa. Ta nhận thấy lựa chọn ``điểm giữa'' thật ra cũng chỉ có \(O(RC)\)
cách. Vậy tại sao không liệt kê điểm giữa, rồi tìm kiếm từ giữa ra hai đầu?

Do đó ta liệt kê điểm giữa của đường đi, gọi là \(\operatorname{Mid}\), tức ô
thứ \(11\). Từ \(\operatorname{Mid}\) tìm kiếm \(10\) bước, ta có thể thu
được một đường đi \(P\) gồm \(11\) ô. Gọi \(\operatorname{Num}(P)\) là tập các
giá trị của \(10\) ô không tính \(\operatorname{Mid}\). Gọi giá trị trong
\(\operatorname{Mid}\) là \(X\). Nếu tồn tại hai đường đi \(P_1\) và \(P_2\)
thỏa mãn
\[
  \operatorname{Num}(P_1)\cup\operatorname{Num}(P_2)\cup\{X\}
  =
  \{0,1,\ldots,20\},
\]
thì \(P_1\) và \(P_2\) ghép lại là một đường đi thỏa mãn yêu cầu. Chú ý rằng
vì các số \(0\) tới \(20\) đều xuất hiện đúng một lần, nên có thể khẳng định
\(P_1\) và \(P_2\) chỉ giao nhau tại \(\operatorname{Mid}\). Giống bài trước,
ta dùng một số nhị phân để biểu diễn tập \(\operatorname{Num}(P)\), nhờ đó có
thể nhanh chóng kiểm tra
\[
  \operatorname{Num}(P_1)\cup\operatorname{Num}(P_2)\cup\{X\}
  =
  \{0,1,\ldots,20\}.
\]

Ước lượng lượng tính toán của thuật toán trên. Trước hết số cách liệt kê
\(\operatorname{Mid}\) là \(O(RC)\). Tiếp theo, vì tìm kiếm \(10\) bước và
mỗi bước có thể đi theo bốn hướng, số đường đi là \(4^{10}\). Thực ra ước
lượng này có thể chính xác hơn: từ bước tìm kiếm thứ hai trở đi, ``đi ngược
lại'' hiển nhiên không hợp lệ, nên từ bước thứ hai chỉ còn ba hướng có thể
đi. Vì vậy số đường đi thực tế không vượt quá \(4\cdot 3^9\). Với mỗi đường
đi, thao tác bảng băm và so sánh đều chỉ cần thời gian hằng số. Tóm lại, ta
thu được thuật toán có độ phức tạp
\[
  O(RC\cdot 4\cdot 3^9).
\]
Dù \(RC\cdot 4\cdot 3^9\) có thể đạt
\[
  21^2\cdot 4\cdot 3^9=34720812,
\]
nhưng do đường đi không được vượt khỏi biên ma trận hoặc đi tới ô không có
số, và nếu một số xuất hiện hai lần trên đường đi thì có thể cắt nhánh ngay,
lượng tìm kiếm thực tế sẽ nhỏ hơn rất nhiều so với ước lượng trên. Trong kiểm
thử thực tế, chương trình cài đặt theo thuật toán trên vượt qua toàn bộ dữ
liệu kiểm thử.

Ví dụ 5 cho thấy khi dùng kỹ thuật \emph{meet in the middle}, nếu điểm đầu
và điểm cuối chưa biết, đồng thời số điểm giữa khá ít, ta có thể thử đổi
``gặp nhau ở giữa'' thành ``xuất phát từ giữa''.

\subsubsection{Ví dụ 6: \emph{FencingGarden}}

\paragraph{Nội dung bài toán.}
Có \(N\) thanh gỗ, độ dài lần lượt là \(L_1,L_2,\ldots,L_N\). Bạn có thể chọn
trước một thanh gỗ và cắt nó thành hai đoạn, độ dài không nhất thiết là số
nguyên. Từ các thanh gỗ thu được, chọn một phần để tựa vào một bức tường dài
vô hạn và rào một vùng hình chữ nhật. Hãy tìm độ dài cạnh song song với tường
của hình chữ nhật khi diện tích vùng hình chữ nhật đạt giá trị lớn nhất. Bài
gốc là TopCoder Single Round Match 461, Division I, Level Three.

\paragraph{Quy mô dữ liệu.}
\[
  N\le 40,\qquad L_i\le 10^8.
\]

\paragraph{Phân tích.}
Bài này khác mọi ví dụ trước: nó không yêu cầu đếm số phương án, mà yêu cầu
một nghiệm tối ưu.

Một khó khăn của bài là có rất nhiều cách cắt thanh gỗ; liệt kê cách cắt hiển
nhiên là không thực tế.

Vì ta cắt một thanh gỗ thành hai đoạn, còn hình chữ nhật được rào bởi ba
cạnh, chắc chắn có một cạnh hoàn toàn được ghép từ các thanh gỗ nguyên vẹn.
Cạnh này có hai khả năng: song song với tường hoặc vuông góc với tường. Gọi
độ dài cạnh đó là \(\operatorname{Len}\).

Trước hết xét trường hợp cạnh này song song với tường. Đặt
\[
  S=\sum_{i=1}^{N} L_i.
\]
Khi đó hai cạnh còn lại của hình chữ nhật phải có độ dài bằng nhau. Ta xếp
các thanh gỗ chưa dùng lại với nhau; tổng độ dài của chúng là
\(S-\operatorname{Len}\). Cắt một nhát tại vị trí
\((S-\operatorname{Len})/2\). Nếu vị trí này đúng là chỗ nối của hai thanh gỗ,
ta không làm thay đổi số lượng thanh gỗ; nếu đây là một điểm ở giữa một thanh
gỗ, ta cắt thanh đó thành hai đoạn. Dù thế nào, ta luôn có thể chia các thanh
gỗ còn lại thành hai phần có độ dài đều là
\((S-\operatorname{Len})/2\), nhờ đó thu được hình chữ nhật diện tích
\[
  \operatorname{Len}\cdot\frac{S-\operatorname{Len}}{2}.
\]

Tương tự, nếu một cạnh độ dài \(\operatorname{Len}\) vuông góc với tường và
hoàn toàn được ghép từ các thanh gỗ nguyên vẹn, ta có thể thu được một hình
chữ nhật diện tích
\[
  \operatorname{Len}(S-2\operatorname{Len}).
\]

Theo tính chất của hàm bậc hai, trong trường hợp thứ nhất
\(\operatorname{Len}=S/2\) là tối ưu; trong trường hợp thứ hai
\(\operatorname{Len}=S/4\) là tối ưu. Vì vậy ta cần tìm, trong các độ dài có
thể ghép từ các thanh gỗ nguyên vẹn, những độ dài gần \(S/2\) và \(S/4\) nhất.

Giống các bài trước, ta vẫn chia \(N\) thanh gỗ thành hai phần, lần lượt có
\(M\) thanh và \(N-M\) thanh. Ta tìm kiếm trước \(M\) thanh và ghi lại các độ
dài thu được. Sau đó tìm kiếm \(N-M\) thanh còn lại; giả sử thu được một phần
có độ dài \(T\), thì cần tìm trong các độ dài đã tìm trước đó một \(x\) sao
cho \(x+T\) gần \(S/2\) hoặc \(S/4\) nhất, tức làm cho \(x\) gần
\(S/2-T\) hoặc \(S/4-T\) nhất. Một cách xử lý tốt là sắp xếp trước
\(O(2^M)\) độ dài thu được ở phần thứ nhất; khi truy vấn ở phần thứ hai, dùng
hai lần tìm kiếm nhị phân.

Sắp xếp \(O(2^M)\) số có độ phức tạp thời gian \(O(2^M M)\). Thực hiện
\(O(2^{N-M})\) lần tìm kiếm nhị phân trong đó có độ phức tạp
\(O(2^{N-M}M)\). Nếu chọn \(M=\lfloor N/2\rfloor\), độ phức tạp thời gian của
thuật toán là
\[
  O(2^{N/2}N).
\]

Ví dụ 6 cho thấy kỹ thuật \emph{meet in the middle} không chỉ giải được các
bài toán đếm, mà còn có thể kết hợp với tìm kiếm nhị phân và các phương pháp
khác để xử lý một số bài tối ưu đặc biệt.

\subsection{Tổng kết}

Từ các phần trên có thể thấy, \emph{meet in the middle} là một kỹ thuật tối
ưu hóa tìm kiếm rất hiệu quả. Tư tưởng cốt lõi của nó là chia bài toán thành
hai phần để tìm kiếm riêng, rồi nhanh chóng hợp nhất kết quả của hai phần tìm
kiếm nhằm tối ưu thuật toán. Tuy ta không nhờ đó thu được một thuật toán đa
thức, nhưng thường có thể làm số mũ trong độ phức tạp thời gian giảm một nửa;
trong một số bài, điều này đã đủ để ta tính được kết quả trong giới hạn thời
gian.

Ưu điểm của kỹ thuật này là thuật toán dễ cài đặt; quá trình tìm kiếm mà
thuật toán dùng về cơ bản giống thuật toán thô, chỉ cần thêm một số cấu trúc
dữ liệu khá đơn giản, chẳng hạn bảng băm. Điểm yếu là phạm vi sử dụng không
quá rộng, và độ phức tạp thời gian sau tối ưu vẫn là cấp số mũ, nên quy mô dữ
liệu có thể xử lý hiệu quả không tăng quá nhiều. Đồng thời, thuật toán phải
tiêu tốn rất nhiều không gian.

Mô hình đồ thị có hướng và mô hình phương trình khái quát các đặc trưng phổ
biến của những bài mà kỹ thuật \emph{meet in the middle} có thể giải; dùng
hai mô hình này có thể giải được một lớp bài toán. Mặt khác, vẫn có những bài
không thể chuyển trực tiếp thành các mô hình trên, đòi hỏi ta sửa kỹ thuật
\emph{meet in the middle}. Từ đó có thể thấy cùng một kỹ thuật cũng có nhiều
cách sử dụng khác nhau trong các bài khác nhau. Ta cần linh hoạt dùng kỹ
thuật \emph{meet in the middle} theo tình huống cụ thể để đạt mục đích tối ưu
thuật toán.

\subsection{Lời cảm ơn}

Xin cảm ơn người thân và bạn bè đã động viên, ủng hộ tôi.

Xin cảm ơn các thầy cô Li Shu, Wu Xiaoshi và Shi Zhenlei đã hướng dẫn tôi.

Xin cảm ơn Jia Zhipeng, Xu Haoran và nhiều bạn học thi tin học khác đã giúp
đỡ và gợi mở cho tôi.

\subsection{Tài liệu tham khảo}

\begin{enumerate}
  \item Thomas H. Cormen, Charles E. Leiserson, Ronald L. Rivest, Clifford
  Stein. \emph{Introduction to Algorithms}.
  \item vexorian. TopCoder Algorithm Problem Set Analysis.
  \item Bruce Merry. USACO 2012 US Open Gold Division, Balanced Cow Subsets
  Solution.
  \item Zhipeng Jia. TopCoder SRM 450--499 Solution.
\end{enumerate}

\end{document}
